%% 第31回編集委員会（物理第5巻）：組 原W3
%% p5-06.tex  原-6　壁をすりぬける：トンネル効果と調和振動子（★3〜5・3段）
%% 中身は計画書 tools/phys5-plan.md の B「原-6」。公式の持ち主は A（A50・A51・A56）。
%% 3段：目でわかる（①②、★3）→ 手でわかる（③〜⑤、★3。ここまでで十分）→ 式でわかる（⑥〜⑫、★5）。
%%   手でわかるの段は、κ の式と e^{−2κa} を「結果を読む」形で使い、電卓（5-14 の e）だけでたどれるように書いた。
%%   有限の井戸・正確な透過率・調和振動子は★5（計画書 B のはじめ）。
%% 中学・高校の公式はこの章にない（N-10(c)）ので \phkoushiki は置かない。
%% 数値はすべて Python で計算した：checks/物理第5巻/原-6/calc_p5_06.py（出力 calc_p5_06_out.txt）、
%%   fig_p5_06.py（⑦の波束が壁に当たる図：クランク・ニコルソン法。30 fs 後にぬけた確率 0.305、平面波の透過率 0.303）。
%% 記号（原-5 と同じ ψ・U・E・i。計画書 C-1）：壁の高さ U₀、壁の厚さ a（⑥の井戸では幅）、減衰の定数 κ、波数 k、
%%   透過率 𝒯・反射率 ℛ（温度 T・気体定数 R とぶつけないため筆記体。E の4(g)：組の判断）、
%%   ばね定数 k_s（熱-15 と同じ。波数 k と分ける）、角振動数 ω、零点のゆれの幅 x₀、換算質量 μ（原-7 ⑪でも使う）、⑥の井戸の係数 c₁・c₂（⑪の行き来の強さ A と分ける。第32回 F3-31）、
%%   分光の波数 ν̃（cm⁻¹）、アインシュタイン温度 θ_E（熱-15 と同じ）、温度 T、気体定数 R。
%% 前の巻・数学との約束（計画書 A-6・C-6-2）：
%%   Q4（力-8：原子のばね・ヤング率）→ ⑨・例題7（CO のばね定数 1.9×10³ N/m）。
%%   Q10（熱-14 ④：分子の振動がふだんの温度でエネルギーを受けとらない）→ ⑩・例題8（CO θ_E＝3080 K、C/R＝0.0036）。
%%   Q11（熱-15 ⑧：Δε は ω に比例。比例の係数は原-6）→ ⑧（Δε＝ħω）。Q14（熱-15：いちばん低いのは ½Δε）→ ⑧（零点エネルギー）。
%%   Q28（波-3：分子の振動数の測定は物質を調べる道具）→ ⑨・例題9（HCl 3.47 μm）。
%%   波-7 ④の反射係数 r＝(Z₁−Z₂)/(Z₁+Z₂)、通りぬける割合 4Z₁Z₂/(Z₁+Z₂)² → ②⑦（A56 の相似。Z を k に）。
%%   波-12 ④の全反射 → ①（しみ出す波は名前だけ）。原-4 ⑫のメーザー → ⑪（アンモニア）。原-12 ⑩のα崩壊 → ⑤で予告。
%% 出典（11月11日に確かめた）
%%   ・定数：CODATA 2022（NIST）。ħ²/(2mₑ)＝0.0381 eV·nm²、κ＝5.12√(U₀−E〔eV〕) nm⁻¹。原子の質量：NIST（AME2020）。
%%   ・分子の定数：NIST Chemistry WebBook（webbook.nist.gov、Huber & Herzberg の2原子分子の定数）を curl で読んだ。
%%       CO（ID=C630080）：ωe＝2169.81358、ωexe＝13.28831 cm⁻¹、re＝1.128323 Å → 0→1 の吸収 2143.24 cm⁻¹。
%%       HCl（C7647010、H³⁵Cl）：ωe＝2990.946、ωexe＝52.8186 cm⁻¹、re＝1.27455 Å → 2885.31 cm⁻¹。
%%       N₂（C7727379）：ωe＝2358.57、ωexe＝14.324 cm⁻¹ → 2329.92 cm⁻¹。
%%   ・ケイ素と酸化ケイ素の境目の電子の壁の高さ 約3.1〜3.2 eV（光電子放出の測定 3.1 eV、ほかに 3.14・3.20 eV）、
%%       酸化膜の中の電子の有効質量 約0.42 mₑ：WebSearch の抜粋で arXiv:1604.03353 ほか複数の論文を照らし合わせた。本文は 3.2 eV。
%%   ・フラッシュメモリの酸化膜の厚さ（約 8〜10 nm、10年の保持をめざす）、書きこむときの電場 約10 MV/cm：
%%       米国特許の公開文書（US 7154153 ほか）と Microchip の SuperFlash の技術資料の抜粋（WebSearch）で照らし合わせた。
%%       「1秒に壁に当たる回数 10¹³」は例題のために置いた仮の値（本文で「〜とします」と書いた）。
%%   ・STM の仕事関数 約4.5 eV（タングステン・金などの金属の典型値。計画書の〔要〕のまま、桁の目安として使う）。
%%   ・アンモニアの反転の線（J＝3、K＝3）23 870 MHz：NBS（米国標準局）の論文 tf.nist.gov/general/pdf/208.pdf と、
%%       京都大学の分光器の論文（WebSearch の抜粋）で照らし合わせた。メーザー（1954年、ゴードン・ザイガー・タウンズ）。
%%   ・ノーベル物理学賞：1964年（タウンズ・バソフ・プロホロフ）、1973年（江崎玲於奈・ジエーバー／ジョセフソン）、
%%       1986年（ビーニッヒ・ローラー：STM）、2025年（クラーク・デヴォレ・マーティニス「電気回路での巨視的な量子トンネルと
%%       エネルギーの量子化の発見」）：nobelprize.org（計画書 D28・D30・D31・D39）。
%%   ・トンネルにかかる時間：U. S. Sainadh ほか, Nature 568, 75 (2019)、arXiv:1707.05445（原子の水素、遅れは 1.8 アト秒以下）。
%%       R. Ramos ほか, Nature 583, 529 (2020)（冷やしたルビジウムの原子、スピンの回転を時計にした）。
%%       2つの時計がちがうものをはかっているという議論：arXiv:2503.07859（2025年）。いずれも WebSearch の抜粋と arXiv の要旨。
%%   ・STM で原子を並べた（35個のキセノン、ニッケルの表面、約4 K）：D. M. Eigler, E. K. Schweizer, Nature 344, 524 (1990)〔二次資料で照合〕。
%% 〔要確認〕：STM の仕事関数 4.5 eV（理科年表に届かず。金属の典型値の目安として「約」をつけた）。
%%   出典の組（11月11日）：今回届くようになった一次資料（IAEA・PDG・NCNR・ICRP など）に仕事関数の表はなく、〔要確認〕のまま。
\makeatletter
\tikzset{txg lab/.style={font=\sffamily\footnotesize, text=bkInk},
  txg wall/.style={fill=bkGray!25, draw=bkInk, line width=0.6pt},
  txg psi/.style={bkBlue, line width=1.2pt},
  txg psi2/.style={bkOrange!85!black, line width=1.2pt, dash pattern=on 3pt off 1.5pt},
  txg pot/.style={bkInk, line width=1.0pt},
  txg lev/.style={bkInk!70, line width=0.6pt, dashed},
  txg arr/.style={bkNavy, line width=0.8pt, -{Stealth[length=1.8mm, width=1.4mm]}},
  txg dim/.style={bkGray, line width=0.5pt, {Stealth[length=1.6mm]}-{Stealth[length=1.6mm]}}}
\newcommand\txgel[1]{\node[circle, fill=bkBlue!80!black, inner sep=0pt, minimum size=2.6mm, font=\sffamily\bfseries\footnotesize, text=white] at #1 {$-$};}
\makeatother

\phchapter{6}{壁をすりぬける：\\トンネル効果と調和振動子}{USBメモリは、どうやって電子を閉じこめて記録する？}{\bklevelrange{3}{5}}{壁をすりぬける：トンネル効果と調和振動子}
\bkmichishirube{5-14,6-2,6-6,6-11,7-8,波-7,波-12,熱-14,熱-15,原-5}{原-7,原-8,原-9,原-12}

\begin{goal}
\begin{itemize}[label={\color{bkTheme}●}, leftmargin=1.6\zw, itemsep=0pt, topsep=1pt]
\item \bklevel{3}\enspace 壁の中で波動関数が $e^{-\kappa x}$ で弱まることを使い、すりぬける割合 $e^{-2\kappa a}$ を見積もれる。
\item \bklevel{3}\enspace フラッシュメモリや走査トンネル顕微鏡のしくみを、トンネル効果で説明できる。
\item \bklevel{5}\enspace 境目のつながりから、階段の反射率・有限の井戸の準位・正確な透過率を求められる。
\item \bklevel{5}\enspace 調和振動子の準位 $(n+\tfrac12)\hbar\omega$ を使い、分子の振動の光と比熱を説明できる。
\end{itemize}
\end{goal}

\begin{hatena}
\noindent\begin{minipage}[c]{0.52\linewidth}\raggedright
写真や音楽を保存する USB メモリや SSD は、電池をつながなくても、何年も記録を忘れません。中では、1つ1つの記録の場所に電子を閉じこめて、「電子がある・ない」で 1 と 0 を記録しています。

電子は、どうやって閉じこめられ、何年もとどまっているのでしょう。そして、書きかえるときは、どうやって出し入れするのでしょう。
\end{minipage}\hfill
\begin{minipage}[c]{0.45\linewidth}\centering
\begin{tikzpicture}[x=0.72mm, y=1mm]
  \fill[bkGreen!15] (0,0) rectangle (56,6);
  \node[txg lab] at (28,2.4) {ケイ素};
  \fill[bkYellow!30, draw=bkInk, line width=0.4pt] (10,6) rectangle (46,8.2);
  \fill[bkGray!35, draw=bkInk, line width=0.5pt] (10,8.2) rectangle (46,14);
  \fill[bkYellow!30, draw=bkInk, line width=0.4pt] (10,14) rectangle (46,17);
  \fill[bkPurple!25, draw=bkInk, line width=0.5pt] (10,17) rectangle (46,23);
  \txgel{(17,11.1)} \txgel{(23,11.1)} \txgel{(29,11.1)} \txgel{(35,11.1)}
  \node[txg lab, anchor=west] at (47,20) {ゲート};
  \node[txg lab, anchor=west] at (47,11.1) {電子の部屋};
  \draw[bkInk, line width=0.4pt] (46.5,7.1) -- (50,4.5);
  \node[txg lab, anchor=north west] at (49,4.6) {うすいガラス};
\end{tikzpicture}
\end{minipage}
\end{hatena}

\noindent 原-5 では、電子を「限りなく高い壁」の箱に閉じこめました。でも、本物の壁の高さは限りがあります。高さに限りのある壁の所で、電子の波はどうなるでしょう。この章では、電子の波が壁の中へしみこみ、うすい壁なら向こう側へ通りぬけてしまう\textbf{トンネル効果}を学びます。後半では、ばねのような位置エネルギーの中の電子や原子（\textbf{調和振動子}）を解き、物理シリーズの『熱学』（熱-15）で使った「エネルギーの粒」の大きさが $\hbar\omega$ であることを確かめます。

\begin{yattemiyou}[コップの水で、光が通りぬけるのを見る]
\textbf{用意するもの}：水を入れたガラスのコップ（側面がまっすぐなもの）、ばね（またはゴムひも）、おもり2つ\par\vspace{2pt}
\noindent\begin{minipage}[c]{0.62\linewidth}
\begin{enumerate}[label={\bkcircled[bkGreen]{\arabic*}}, leftmargin=2.4\zw, itemsep=1pt, topsep=0pt]
\item コップに水を入れ、指でコップの側面を持つ。上から水面ごしに、ななめに側面を見下ろす。側面が鏡のように光って見える角度をさがす（全反射：波-12 ④）。
\item その角度のまま、指を側面に強く押し当てる。指の当たった所に、指紋の線だけがくっきり見える。
\item ばねにおもりを1つつるしてゆらし、10回ゆれる時間をはかる。ばねを2本並べて（ばね定数を2倍にして）同じようにはかる。
\end{enumerate}
\end{minipage}\hfill
\begin{minipage}[c]{0.35\linewidth}\centering
\begin{tikzpicture}[x=1mm, y=1mm]
  \draw[bkInk, line width=0.7pt] (0,0) -- (0,30) (22,0) -- (22,30) (0,0) -- (22,0);
  \fill[bkBlue!12] (0.4,0.4) rectangle (21.6,24);
  \draw[bkBlue!60, line width=0.5pt] (0.4,24) -- (21.6,24);
  \fill[bkOrange!35, draw=bkInk, line width=0.4pt] (22,8) rectangle (27,15);
  \draw[bkRed!80!black, line width=0.8pt, -{Stealth[length=1.6mm]}] (14.4,2) -- (21.6,13) -- (14.4,24) -- (7.4,31);
  \node[txg lab, anchor=south east] at (7.4,31) {目};
  \node[txg lab, anchor=west] at (27.5,11.5) {指};
\end{tikzpicture}
\end{minipage}
\end{yattemiyou}

\noindent やってみようの答え：①②全反射している所に指を押し当てると、指紋の山（皮ふがガラスにぴったり触れた所）だけ、光が全反射せずに指の中へ通りぬけ、暗い線になって見えます。指紋の谷は、ガラスとの間に空気のすき間があるので、全反射のままです。すき間が光の波長（約 \qty{0.5}{\micro m}）よりずっと広いと、光は通りぬけられません。③ばね定数を2倍にすると、10回ゆれる時間は約 $1/\sqrt2\approx0.71$ 倍になり、振動数は約1.4倍になります（6-6・波-1。$\omega=\sqrt{k_{\mathrm s}/m}$）。

\begin{omoidasou}
\bkrecall{原-5}{電子の波は $-\dfrac{\hbar^2}{2m}\psi''+U\psi=E\psi$ を満たす（⑥）。$E-U$ が大きいほど強く曲がる。$|\psi|^2$ が見つかる確率（④）。}
\bkrecall{6-11}{$y''=a^2y$ の解は $e^{ax}$・$e^{-ax}$（$\cosh ax$・$\sinh ax$）。$\sinh x=(e^x-e^{-x})/2$（②⑤）。}
\bkrecall{波-7}{2本のひもの境目の反射係数は $r=\dfrac{Z_1-Z_2}{Z_1+Z_2}$、通りぬける割合は $\dfrac{4Z_1Z_2}{(Z_1+Z_2)^2}$（④）。}
\bkrecall{熱-15}{ゆれる原子のエネルギーは粒 $\Delta\varepsilon$ の何個分か。$\theta_{\mathrm E}=\Delta\varepsilon/k_{\mathrm B}$ より低い温度で比熱が小さくなる（⑧）。}
\end{omoidasou}

\noindent この章の記号：壁の高さを $U_0$、壁の厚さを $a$ と書きます。壁を通りぬける割合（\textbf{透過率}）は筆記体の $\mathcal T$（ティー）、はね返る割合（\textbf{反射率}）は $\mathcal R$（アール）と書きます（温度 $T$・気体定数 $R$ と分けるため）。ばね定数は $k_{\mathrm s}$ と書き、波数 $k$ と分けます。

%% ======================== 物理の道具 ========================
\bksection{doguu}{物理の道具}

\phstage{me}{壁の所で、波はどうなる？}{3}

\bksubsection{1}{坂をこえられないボールと、しみこむ波}
\label{ph:原-6:sec1}
ボールを転がして、坂を登らせます。ボールの運動エネルギーが、坂の上の位置エネルギーより小さければ、ボールは坂の途中で止まって、もどってきます。坂の向こうに出ることは、けっしてありません（物理シリーズの『力学』（力-4）のエネルギー保存）。

\noindent 電子の波ではどうでしょう。原-5 ⑥の読み方で、「$E-U$ が小さい所ほど、波はゆるく曲がる」と見ました。$E<U$ の所（ボールなら入れない所）では、$E-U$ が負になり、波の曲がる向きが逆になります。すると、波はゆれなくなり、奥へ行くほど小さくなっていきます。これを「波が壁の中へ\textbf{しみこむ}」と言います（図\ref{ph:原-6:fig-shimi}）。

\begin{bkfigure}
\begin{tikzpicture}[x=1mm, y=1mm]
  \begin{scope}
    \fill[bkGray!20] (30,0) -- (42,14) -- (54,14) -- (54,0) -- cycle;
    \draw[txg pot] (0,0) -- (30,0) -- (42,14) -- (54,14);
    \shade[ball color=bkOrange!70] (12,3) circle (3);
    \draw[txg arr] (17,3) -- (27,3);
    \draw[txg arr] (31,9) -- (19,9);
    \node[txg lab, anchor=south] at (25,9.5) {もどる};
    \node[txg lab, anchor=north] at (27,-2) {(ア) ボールは、坂の途中でもどる};
  \end{scope}
  \begin{scope}[shift={(66,0)}]
    \fill[bkGray!20] (30,0) rectangle (40,22);
    \draw[txg pot] (0,0) -- (30,0) -- (30,22) -- (40,22) -- (40,0) -- (66,0);
    \draw[bkRed!80!black, line width=0.8pt, dashed] (0,14) -- (66,14);
    \node[txg lab, anchor=south west] at (0,14.5) {電子のエネルギー $E$};
    \node[txg lab, anchor=south] at (35,22.5) {壁 $U_0$};
    % 境目で値と傾きがつながるように：左右の sin の位相 φ は tan φ = −k/κ（k＝1.2、κ＝0.12）で φ＝95.7°
    \draw[txg psi] plot[domain=0:30, samples=80] (\x, {7+3.2*sin(deg(1.2*(\x-30)) + 95.7)});
    \draw[txg psi] plot[domain=30:40, samples=40] (\x, {7+3.2*sin(95.7)*exp(-0.12*(\x-30))});
    \draw[txg psi] plot[domain=40:66, samples=80] (\x, {7+3.2*exp(-1.2)*sin(deg(1.2*(\x-40))+95.7)});
    \draw[bkGray, line width=0.4pt] (0,7) -- (66,7);
    \node[txg lab, anchor=north] at (33,-2) {(イ) 電子の波は、壁の中で弱まり、向こうへ少し出る};
  \end{scope}
\end{tikzpicture}
\bkfigcaption{(ア) ボールと坂。(イ) 電子の波と、高さ $U_0$・厚さに限りのある壁（灰色）。波は壁の中で、ゆれずに小さくなる。壁がうすいと、向こう側に小さな波が残る。細い線は波の0の位置。}\label{ph:原-6:fig-shimi}
\end{bkfigure}

\noindent 壁が厚ければ、波は壁の中で0近くまで小さくなり、全部はね返ります。でも、壁がうすいと、しみこんだ波が壁の向こう側にとどき、そこからまた波として進みます。つまり、電子が、こえられないはずの壁の向こうで、ある確率で見つかるのです。山にトンネルをほって通りぬけるようなので、これを\textbf{トンネル効果}といいます。\bkidx{とんねるこうか}{トンネル効果}

\begin{phmijika}[光も、全反射の境目の向こうへしみ出す]
物理シリーズの『波動』（波-12 ④）の全反射では、光は境目で全部はね返りました。じつは、全反射している光も、境目の向こうの空気の中へ、波長ほどの深さだけしみ出しています（専門の本では「エバネッセント波」とよぶ）。そこへ別のガラスや指を波長より近づけると、光の一部が通りぬけます。やってみようの指紋が見えたのは、この光の通りぬけの、すき間が0になった場合です。光の「トンネル効果」は、光ファイバーの光を分ける部品などに使われています。
\end{phmijika}

%% ---------------------------------------------------------------
\bksubsection{2}{階段での反射：波-7 と同じ形}
\label{ph:原-6:sec2}
こんどは、電子のエネルギー $E$ が、段差の高さ $U$ より大きい場合を考えます（図\ref{ph:原-6:fig-kaidan}）。ボールなら、段の上でゆっくりになるだけで、必ず全部が段を上がります。ところが電子の波では、段の所で、一部がはね返るのです。

\noindent 段の左と右で、電子の運動エネルギーは $E$ と $E-U$ です。波数は、運動エネルギーを $K$ とすると $k=\sqrt{2mK}/\hbar$ です（原-5 ⑥）。
\[k_1=\frac{\sqrt{2mE}}{\hbar},\qquad k_2=\frac{\sqrt{2m(E-U)}}{\hbar}\]
はね返る割合（反射率）は、⑦で導くように
\pheq{\mathcal R=\Bigl(\frac{k_1-k_2}{k_1+k_2}\Bigr)^2}{単位なし}
\label{ph:原-6:hansha}です。物理シリーズの『波動』（波-7 ④）の、2本のひもの境目の反射係数
\[r=\frac{Z_1-Z_2}{Z_1+Z_2}\]
を2乗したものと、同じ形です。ひものインピーダンス $Z$ の役を、電子では波数 $k$（運動量 $\hbar k$）がしています。違う物理でも、「波が、性質の違う所に入るとき、一部がはね返る」という同じしくみなのです。\bkidx{はんしゃりつ}{反射率（電子の波）}

\begin{bkfigure}
\begin{tikzpicture}[x=0.9mm, y=1mm]
  \fill[bkGray!20] (40,0) rectangle (70,8);
  \draw[txg pot] (0,0) -- (40,0) -- (40,8) -- (70,8);
  \draw[bkRed!80!black, line width=0.8pt, dashed] (0,16) -- (70,16);
  \node[txg lab, anchor=south west] at (0,16.5) {電子のエネルギー $E$};
  \node[txg lab, anchor=south] at (55,8.5) {段 $U$};
  \draw[txg arr] (6,24) -- (30,24); \node[txg lab, anchor=south] at (18,24.5) {来る波 $k_1$};
  \draw[txg arr, dashed] (30,31) -- (6,31); \node[txg lab, anchor=south] at (18,31.5) {はね返る波};
  \draw[txg arr] (44,24) -- (68,24); \node[txg lab, anchor=south, align=center] at (56,24.5) {進む波 $k_2$\\（波長が長い）};
\end{tikzpicture}\hspace{4mm}%
\begin{tikzpicture}
  \begin{axis}[bkplot, width=50mm, height=40mm, xmin=0.8, xmax=6.3, ymin=0, ymax=1.1, xtick={1,2,4,6}, ytick={0.5,1}, xlabel={$E/U$}, ylabel={$\mathcal R$}]
    \addplot[bkcurve, domain=1:6, samples=120] {((sqrt(x)-sqrt(x-1))/(sqrt(x)+sqrt(x-1)))^2};
    \fill[bkRed] (axis cs:2,0.0294) circle (1.8pt);
  \end{axis}
\end{tikzpicture}
\bkfigcaption{段差での電子の波（左）と、反射率 $\mathcal R$ と $E/U$ の関係（右）。$E$ が $U$ に近いほどよくはね返り、$E$ が大きいとほとんど通りぬける。●は例題2。}\label{ph:原-6:fig-kaidan}
\end{bkfigure}

\begin{phbaisuru}[$U$ の近く・10倍]{エネルギー $E$}
$E$ が $U$ に近づくと、段の上の運動エネルギー $E-U$ が0に近づき、$k_2\to0$ です。すると $\mathcal R\to\bigl(k_1/k_1\bigr)^2=1$ で、ほとんど全部はね返ります。逆に $E$ が $U$ の10倍なら、$k_2/k_1=\sqrt{0.9}\approx0.95$ で、$\mathcal R\approx0.0007$、ほとんど全部通りぬけます。エネルギーが大きいと、段差は「小さな段」にしか見えなくなるのです。
\end{phbaisuru}

\Needspace{8\baselineskip}\phstage{te}{すりぬける割合を見積もる}{3}

\bksubsection{3}{壁の中では $e^{-\kappa x}$ で弱まる}
\label{ph:原-6:sec3}
高さ $U_0$ の壁の中（$E<U_0$）で、電子の波は、
\pheq{\psi\propto e^{-\kappa x},\qquad \kappa=\frac{\sqrt{2m(U_0-E)}}{\hbar}}{1/m ＝ $\sqrt{\text{kg·J}}$ ÷ (J·s)}
\label{ph:原-6:kappa}の形で弱まります（$x$ は壁に入ってからの深さ。⑥で導く）。$\kappa$（カッパ）は弱まり方の速さで、波数 $k$ の式の $E-U$ を $U_0-E$ にかえたものです。\bkidx{げんすいのていすう}{減衰の定数 $\kappa$}電子では、長さを nm、エネルギーを eV で書くと、
\[\kappa\approx5.12\times\sqrt{U_0-E\ \text{〔eV〕}}\ \ \text{nm}^{-1}\]
です（$\sqrt{2m\times\qty{1}{eV}}/\hbar=\qty{5.12}{nm^{-1}}$）。壁の高さとの差が \qty{1}{eV}なら、深さ $1/\kappa\approx\qty{0.2}{nm}$ ごとに、波の大きさが $1/e\approx0.37$ 倍になります（5-14）。

\phboxhead{すりぬける割合}
見つかる確率は $|\psi|^2$ なので、厚さ $a$ の壁を通りぬけた所の確率は、入り口の
\[(e^{-\kappa a})^2=e^{-2\kappa a}\]
倍になります。だから、すりぬける割合（透過率）は、だいたい
\pheq{\mathcal T\approx e^{-2\kappa a}}{単位なし}
\label{ph:原-6:toka}です（もっと正確な式は⑦）。\bkidx{とうかりつ}{透過率（トンネル効果）}

\begin{phmodel}[四角い壁のトンネル効果]
\phignore{$e^{-2\kappa a}$ の前につく、1くらいの大きさの係数（⑦）。壁の中の電子の有効質量（原-5 ⑤）。}
\phassume{壁は四角（高さ $U_0$ が厚さ $a$ の間だけ一定）。壁が十分厚い（$\kappa a$ が1より大きい）。}
\phlimit{電圧をかけて壁がななめになるとき（④の書きこみ）。壁が原子1〜2個分よりうすいとき（⑦の正確な式）。}
\end{phmodel}

\begin{phyomi}[$\mathcal T\approx e^{-2\kappa a}$]
「すりぬける割合は、壁の厚さ $a$ に対して、指数関数で急に小さくなる」と読みます。肩の $2\kappa a$ が1ふえるごとに、割合は約 0.37 倍です。壁の高さ（$U_0-E$）は $\kappa$ に $\sqrt{\ }$ で効き、厚さはそのまま効きます。
\end{phyomi}

\begin{bkcenter}
\begin{bktable}{colspec={Q[c,m,wd=26mm] Q[c,m,wd=14mm] *{3}{Q[c,m,wd=24mm]}}, rowsep=1.5pt, row{1}={bg=bkTheme!10, font=\sffamily\small}}
$U_0-E$〔eV〕 & $\kappa$〔nm$^{-1}$〕 & $a=\qty{0.5}{nm}$ & $a=\qty{1}{nm}$ & $a=\qty{2}{nm}$ \\
1 & 5.12 & $\num{6.0e-3}$ & $\num{3.6e-5}$ & $\num{1.3e-9}$ \\
3.2 & 9.17 & $\num{1.0e-4}$ & $\num{1.1e-8}$ & $\num{1.2e-16}$ \\
\end{bktable}
\bktabcaption{すりぬける割合 $e^{-2\kappa a}$ の目安。厚さが2倍になると、割合は2乗になる。}\label{ph:原-6:tab-toka}
\end{bkcenter}

\begin{phbaisuru}{壁の厚さ $a$}
肩の $2\kappa a$ が2倍になるので、割合は $e^{-4\kappa a}=(e^{-2\kappa a})^2$、つまり\textbf{2乗}になります。$\num{3.6e-5}$ だった割合は、$\num{1.3e-9}$ になります。半分ではなく、とてつもなく小さくなるのです。また、壁の高さの差 $U_0-E$ を4倍にすると、$\kappa$ は $\sqrt4=2$ 倍で、これも割合を2乗にします。
\end{phbaisuru}

%% ---------------------------------------------------------------
\bksubsection{4}{フラッシュメモリと走査トンネル顕微鏡}
\label{ph:原-6:sec4}
\phboxhead{電子の部屋とガラスの壁}
USB メモリや SSD の記録の場所（フラッシュメモリの1つの「セル」）は、トランジスタ（原-10）の上に、周りから切りはなされた小さな導体の部屋（\textbf{浮遊ゲート}）をのせたものです（図\ref{ph:原-6:fig-flash}）。\bkidx{ふゆうげーと}{浮遊ゲート}部屋の上下は、二酸化ケイ素（ガラスの仲間）のうすい膜で囲まれています。この膜は、電子にとって高さ約 \qty{3.2}{eV}の壁です。部屋に電子がたまっているかどうかで、下のトランジスタの電流の流れやすさが変わり、それを「0」と「1」として読み出します。

\begin{bkfigure}
\begin{tikzpicture}[x=1mm, y=1mm]
    \fill[bkGreen!15] (-4,0) rectangle (64,8);
    \fill[bkGray!40] (-4,0) rectangle (12,8) (48,0) rectangle (64,8);
    \node[txg lab] at (4,4) {ソース};
    \node[txg lab] at (56,4) {ドレイン};
    \node[txg lab] at (30,3) {ケイ素};
    \fill[bkYellow!30, draw=bkInk, line width=0.4pt] (12,8) rectangle (48,11);
    \fill[bkGray!35, draw=bkInk, line width=0.5pt] (12,11) rectangle (48,18);
    \fill[bkYellow!30, draw=bkInk, line width=0.4pt] (12,18) rectangle (48,22);
    \fill[bkPurple!25, draw=bkInk, line width=0.5pt] (12,22) rectangle (48,29);
    \txgel{(19,14.5)} \txgel{(26,14.5)} \txgel{(33,14.5)} \txgel{(40,14.5)}
    \node[txg lab, anchor=west] at (66,25.5) {制御ゲート};
    \node[txg lab, anchor=west] at (66,14.5) {浮遊ゲート（電子の部屋）};
    \draw[bkInk, line width=0.4pt] (48.5,25.5) -- (65.5,25.5) (48.5,14.5) -- (65.5,14.5) (48.5,9.5) -- (65.5,9.5) (48.5,20) -- (65.5,20);
    \node[txg lab, anchor=west] at (66,9.5) {うすい酸化膜（約 \qty{8}{nm}）};
    \node[txg lab, anchor=west] at (66,20) {酸化膜};
\end{tikzpicture}
\bkfigcaption{フラッシュメモリの1つのセルの断面（たての大きさを広げて描いた）。浮遊ゲートは、周りを酸化膜で囲まれている。}\label{ph:原-6:fig-flash}
\end{bkfigure}


\label{ph:原-6:hatena}\noindent\bkhatenaref の答えです。酸化膜の厚さを約 \qty{8}{nm}とすると、$\kappa=5.12\times\sqrt{3.2}\approx\qty{9.2}{nm^{-1}}$ で、すりぬける割合は $e^{-2\times9.2\times8}\approx10^{-64}$ です。部屋の中の電子が1秒に $10^{13}$ 回ほど壁に当たるとしても、10年で約 $3\times10^{21}$ 回で、すりぬける見こみは全くありません。だから電子は、電池がなくても何年もとどまるのです（例題1）。

\noindent 書きこむときは、制御ゲートに高い電圧（十数 V）をかけ、酸化膜の中に \qty{1}{nm}あたり約 \qty{1}{V}の強い電場を作ります。すると、位置エネルギーの壁が、図\ref{ph:原-6:fig-kabe}(イ)のようにななめになります。電子のエネルギーの高さで見た壁の厚さは、$3.2\div1\approx\qty{3}{nm}$ ほどにうすくなり、電子がトンネル効果で部屋に入ります。消すときは、逆向きの電場で電子をぬき出します。このように、ななめの壁を通りぬけるトンネル効果を、\textbf{ファウラー・ノルドハイムのトンネル}といいます。\bkidx{ふぁうらーのるどはいむのとんねる}{ファウラー・ノルドハイムのトンネル}

\begin{bkfigure}
\begin{tikzpicture}[x=1.1mm, y=1.25mm]
    \fill[bkGray!20] (14,0) rectangle (28,9.6);
    \draw[txg pot] (0,0) -- (14,0) -- (14,9.6) -- (28,9.6) -- (28,0) -- (42,0);
    \txgel{(5,1.6)} \txgel{(9,1.6)}
    \node[txg lab, anchor=south] at (21,10.1) {壁 \qty{3.2}{eV}};
    \draw[txg dim] (14,-2) -- (28,-2); \node[txg lab, anchor=north] at (21,-2.5) {約 \qty{8}{nm}};
    \node[txg lab, anchor=north] at (7,-1) {部屋};
    \node[txg lab, anchor=north] at (21,-8.5) {(ア) ふだん};
    \begin{scope}[shift={(58,0)}]
      \fill[bkGray!20] (14,0) -- (14,9.6) -- (19.6,0) -- cycle;
      \draw[txg pot] (0,0) -- (14,0) -- (14,9.6) -- (28,-14.4) -- (42,-14.4);
      \draw[bkRed!80!black, line width=0.8pt, dashed] (-2,0) -- (44,0);
      \txgel{(5,1.6)}
      \draw[txg arr, line width=1.1pt] (8,1.6) -- (38,1.6);
      \draw[txg dim] (14,-2) -- (19.6,-2);
      \node[txg lab, anchor=north east] at (19.6,-2.6) {約 \qty{3}{nm}};
      \node[txg lab, anchor=north] at (3,-7) {ケイ素};
      \node[txg lab, anchor=south] at (36,-13.9) {部屋};
      \node[txg lab, anchor=north] at (21,-15.5) {(イ) 書きこむとき（強い電場）};
    \end{scope}
\end{tikzpicture}
\bkfigcaption{酸化膜の壁。(ア) ふだんは、電子は高く厚い壁に閉じこめられる。(イ) 強い電場をかけると、壁がななめになり、電子の高さ（破線）で見た壁の厚さが約 \qty{3}{nm}にうすくなって、ケイ素から部屋へ電子がすりぬける。}\label{ph:原-6:fig-kabe}
\end{bkfigure}

\phboxhead{原子を見る：走査トンネル顕微鏡}
とがった金属の針を、物の表面から \qty{1}{nm}ほどまで近づけ、小さな電圧をかけます。針と表面の間は真空で、電子にとっては高さ（仕事関数、原-2）約 \qty{4.5}{eV}の壁です。それでも、トンネル効果で小さな電流が流れます。電流は $e^{-2\kappa d}$（$d$ はすき間）に比例し、$\kappa\approx\qty{11}{nm^{-1}}$ なので、針を \qty{0.1}{nm}近づけるだけで、電流は約9倍になります（例題4）。

\noindent 針で表面をなぞりながら、電流が一定になるように針の高さを動かすと、針の高さの記録が、表面の原子1個1個のでこぼこの地図になります。これが\textbf{走査トンネル顕微鏡}（STM）です。\bkidx{そうさとんねるけんびきょう}{走査トンネル顕微鏡}1981年ごろにビーニッヒとローラーが作り、1986年のノーベル物理学賞を受けました。原子より小さい \qty{0.01}{nm}の高さの違いまで見分けられるのは、電流が距離で指数関数的に変わるからです。

\begin{bkfigure}
\begin{tikzpicture}[x=1mm, y=1mm]
  \foreach \i in {0,...,11} {\shade[ball color=bkGreen!50] (8+\i*7,4) circle (3.4);}
  \fill[bkGray!50, draw=bkInk, line width=0.5pt] (43.5,30) -- (46,14.2) -- (48.5,30) -- cycle;
  \shade[ball color=bkGray!60] (46,13.3) circle (1.5);
  \draw[bkOrange!85!black, line width=1.1pt, -{Stealth[length=1.6mm]}] (46,11.7) -- (46,7.9);
  \draw[txg dim] (56,7.4) -- (56,11.8);
  \node[txg lab, anchor=west] at (57,9.6) {すき間 $d$（約 \qty{1}{nm}）};
  \node[txg lab, anchor=west] at (50,24) {針};
  \node[txg lab, anchor=east] at (36,18) {トンネル電流};
  \draw[bkInk, line width=0.4pt] (36.3,17.5) -- (45,10.2);
  \draw[bkPurple, line width=0.9pt, dashed] plot[domain=4:88, samples=120] (\x, {38+1.2*cos(deg(2*pi*(\x-8)/7))});
  \node[txg lab, anchor=west] at (89,38) {針の高さの記録};
\end{tikzpicture}
\bkfigcaption{走査トンネル顕微鏡のしくみ。表面の原子（球）の上を、電流が一定になるように針の高さを変えながらなぞる。針の高さの記録（破線）が、原子の並びの形になる。}\label{ph:原-6:fig-stm}
\end{bkfigure}

%% ---------------------------------------------------------------
\bksubsection{5}{エサキダイオード、α崩壊、太陽}
\label{ph:原-6:sec5}
トンネル効果は、半導体の部品から星の中まで、いろいろな所で働いています。
\begin{itemize}[leftmargin=1.2\zw, itemsep=2pt, topsep=2pt, label={\color{bkTheme}●}]
\item \textbf{エサキダイオード}：1957年、東京通信工業（今のソニー）の江崎（えさき）は、半導体の pn 接合（原-9）の境目の層をとてもうすくすると、電子がトンネル効果で通りぬけ、電圧を上げると逆に電流が減る範囲が現れることを見つけました。半導体の中のトンネル効果を初めて確かめた発見で、1973年のノーベル物理学賞を受けました。\bkidx{えさきだいおーど}{エサキダイオード}
\item \textbf{α崩壊}：ウランなどの重い原子核は、α粒子（ヘリウムの原子核）を放出してこわれます。α粒子は、原子核の中でクーロン力の高い壁に囲まれていて、そのエネルギーでは壁をこえられません。それでも、トンネル効果で少しずつぬけ出します。壁の厚さのわずかな違いが、ぬけ出す割合を何十けたも変えるので、半減期が1秒より短いものから何十億年のものまである理由が説明されます（1928年、ガモフほか。原-12 ⑩で詳しく）。
\item \textbf{太陽}：太陽の中心では、水素の原子核（陽子）どうしが合体する反応（核融合）でエネルギーを出しています。中心の温度でも、陽子の熱運動のエネルギーは、陽子どうしのクーロン力の壁よりずっと低く、壁をこえられるものはほとんどありません。核融合が起きるのは、トンネル効果で壁をすりぬけるからです（原-11 ⑪）。
\end{itemize}

\phkokomade

\Needspace{12\baselineskip}\phstage{shiki}{境目をつないで解く}{5}

\bksubsection{6}{有限の深さの井戸}
\label{ph:原-6:sec6}
原-5 の箱の壁を、高さ $U_0$ の有限の壁にします。井戸の幅を $a$、まん中を $x=0$ とし、井戸の中（$|x|<a/2$）は $U=0$、外は $U=U_0$ です。$0<E<U_0$ の状態（井戸に閉じこめられた状態）をさがします。

\phboxhead{場所ごとの解}
\begin{itemize}[leftmargin=1.2\zw, itemsep=2pt, topsep=2pt, label={\color{bkTheme}●}]
\item 井戸の中：$\psi''=-k^2\psi$（$k=\sqrt{2mE}/\hbar$）なので、$\cos kx$ か $\sin kx$（原-5 ⑧）。
\item 井戸の外：$\psi''=+\kappa^2\psi$（$\kappa=\sqrt{2m(U_0-E)}/\hbar$）なので、$e^{\kappa x}$ と $e^{-\kappa x}$（6-11 ⑤）。遠くで $\psi$ が限りなく大きくなってはいけないので（確率の合計が1にならない）、右の外は $e^{-\kappa x}$、左の外は $e^{\kappa x}$ だけが残ります。③の「壁の中で $e^{-\kappa x}$ で弱まる」は、このことでした。
\end{itemize}
井戸は左右対称なので、解は左右対称（偶：$\cos kx$）か、反対称（奇：$\sin kx$）のどちらかです。

\phboxhead{境目のつながり}
シュレディンガー方程式で $U$ が有限なら、$\psi''$ も有限なので、$\psi$ と $\psi'$（傾き）は境目で切れ目なくつながります。偶の解で、境目 $x=a/2$ の値と傾きをそろえると、
\[c_1\cos\frac{ka}{2}=c_2e^{-\kappa a/2},\qquad -kc_1\sin\frac{ka}{2}=-\kappa c_2e^{-\kappa a/2}\]
です（$c_1$・$c_2$ は係数）。2つめを1つめでわると、$c_1$ と $c_2$ が消えて、
\[k\tan\frac{ka}{2}=\kappa\]
となります。$z=ka/2$、$z_0=\dfrac a2\cdot\dfrac{\sqrt{2mU_0}}{\hbar}$ とおくと、$\kappa a/2=\sqrt{z_0^{\,2}-z^2}$ なので、
\pheq{\text{偶：}z\tan z=\sqrt{z_0^{\,2}-z^2},\qquad \text{奇：}-\frac{z}{\tan z}=\sqrt{z_0^{\,2}-z^2}}{単位なし}
\label{ph:原-6:yugen}です（奇の式も同じようにして出ます）。この式は、式の変形では解けません（このような方程式を超越方程式といいます）。

\phboxhead{グラフの交点で解く}
左辺と右辺のグラフを描き、交わる点をさがします。右辺 $\sqrt{z_0^{\,2}-z^2}$ は半径 $z_0$ の円の4分の1です。深さ \qty{1}{eV}・幅 \qty{1}{nm}の井戸では $z_0=0.5\times5.12\times1\approx2.56$ で、交点は2つ（偶が1つ、奇が1つ）です（図\ref{ph:原-6:fig-kouten}）。交点の $z$ から $E=\dfrac{\hbar^2}{2m}\Bigl(\dfrac{2z}{a}\Bigr)^2$ を求めると、\qty{0.19}{eV}と \qty{0.70}{eV}です（例題5）。表計算で、$z$ を少しずつ変えて両辺の差の符号が変わる所をさがしたり、範囲を半分ずつにせばめたり（二分法）すれば、何けたでも求められます。

\begin{bkfigure}
\begin{tikzpicture}
\begin{axis}[bkplot, width=66mm, height=44mm, xmin=0, xmax=3.45, ymin=0, ymax=3.9, xtick={1.5708,2.5616}, xticklabels={$\pi/2$,$z_0$}, ytick={1,2,3}, xlabel={$z$}, ylabel={}, clip=true]
  \addplot[bkcurve, domain=0:1.2, samples=80] {x*tan(deg(x))};
  \addplot[bkcurve2, dash pattern=on 4pt off 1.5pt, domain=1.62:2.56, samples=80] {-x/tan(deg(x))};
  \addplot[bkthin, line width=1.1pt, domain=0:2.5616, samples=120] {sqrt(max(2.5616^2-x^2,0))};
  \fill[bkRed] (axis cs:1.1188,2.3021) circle (2pt) (axis cs:2.1474,1.3967) circle (2pt);
  \node[txg lab, anchor=west] at (axis cs:0.05,3.35) {偶（実線）};
  \node[txg lab, anchor=south east] at (axis cs:3.43,1.6) {奇（破線）};
\end{axis}
\end{tikzpicture}
\bkfigcaption{深さ \qty{1}{eV}・幅 \qty{1}{nm}の井戸の、グラフの交点（●）。$z\tan z$（実線）と $-z/\tan z$（破線）が、4分の1の円（細い線）と交わる。}\label{ph:原-6:fig-kouten}
\end{bkfigure}

\begin{bkfigure}
\begin{tikzpicture}[x=1mm, y=1mm]
  \fill[bkGray!20] (-34,0) rectangle (-15,40) (15,0) rectangle (34,40);
  \draw[txg pot] (-34,40) -- (-15,40) -- (-15,0) -- (15,0) -- (15,40) -- (34,40);
  \node[txg lab, anchor=south] at (24.5,40.5) {$U_0=\qty{1}{eV}$};
  \draw[txg lev] (-34,7.6) -- (34,7.6);
  \draw[txg lev] (-34,28.1) -- (34,28.1);
  \draw[txg psi] plot[domain=-15:15, samples=60] (\x, {7.6+6*cos(deg(2.2378*\x/30))});
  \draw[txg psi] plot[domain=15:34, samples=40] (\x, {7.6+6*0.4367*exp(-4.6086*(\x-15)/30)});
  \draw[txg psi] plot[domain=-34:-15, samples=40] (\x, {7.6+6*0.4367*exp(4.6086*(\x+15)/30)});
  \draw[txg psi2] plot[domain=-15:15, samples=60] (\x, {28.1+6*sin(deg(4.2949*\x/30))});
  \draw[txg psi2] plot[domain=15:34, samples=40] (\x, {28.1+6*0.8383*exp(-2.7929*(\x-15)/30)});
  \draw[txg psi2] plot[domain=-34:-15, samples=40] (\x, {28.1-6*0.8383*exp(2.7929*(\x+15)/30)});
  \node[txg lab, anchor=west] at (35,7.6) {\qty{0.19}{eV}};
  \node[txg lab, anchor=west] at (35,28.1) {\qty{0.70}{eV}};
  \draw[txg dim] (-15,-3) -- (15,-3); \node[txg lab, anchor=north] at (0,-3.5) {\qty{1}{nm}};
\end{tikzpicture}
\bkfigcaption{深さ \qty{1}{eV}・幅 \qty{1}{nm}の井戸の2つの準位（破線）と波動関数（実線と破線の曲線）。波は壁の中へしみ出している。}\label{ph:原-6:fig-yugen}
\end{bkfigure}

\noindent 有限の井戸の準位は（図\ref{ph:原-6:fig-yugen}）、同じ幅の無限の井戸（\qty{0.376}{eV}・\qty{1.50}{eV}）より低くなります。波が壁の中へしみ出して、実際の幅が少し広がったのと同じだからです。また、$z_0$ がどんなに小さくても、偶の交点は必ず1つあります。1次元の井戸には、どんなに浅くても、閉じこめられた状態が少なくとも1つあるのです。交点の数は、$z_0$ を $\pi/2$ でわって切り上げた数です。

\phboxhead{差分法でも確かめる}
原-5 ⑫の差分法のプログラムで、位置エネルギー \texttt{U} を井戸の形にするだけで、同じ準位が出ます。井戸の外の壁も、遠くまで（ここでは $\pm\qty{3}{nm}$ まで）点を並べて入れておきます。
\begin{phcode}[Python]
import numpy as np
L, N = 3.0, 1200                  \# x は −3〜3 nm、点の数 N
dx = 2 * L / (N + 1)
x = -L + dx * np.arange(1, N + 1)
U = np.where(abs(x) < 0.5, 0.0, 1.0)   \# 井戸の中 0 eV、外 1 eV
a = 0.0381 / dx**2
H = (np.diag(2*a + U) + np.diag(-a*np.ones(N-1), 1)
     + np.diag(-a*np.ones(N-1), -1))
E = np.linalg.eigvalsh(H)         \# 固有値だけを小さい順に
print(E[E < 1].round(3))          \# 壁より低い（閉じこめられた）準位
\end{phcode}
\begin{phcodeout}[実行結果]
[0.191 0.703]
\end{phcodeout}
グラフの交点の答え（0.19・\qty{0.70}{eV}）と合います。\texttt{np.where} は「条件が正しい所は 0.0、そうでない所は 1.0」という配列を作ります。\texttt{U} の式を変えれば、2つの井戸（⑪）や、ななめの壁（④）でも、同じ数行で準位が求まります。

%% ---------------------------------------------------------------
\bksubsection{7}{境目をつないで、反射率と透過率を導く}
\label{ph:原-6:sec7}
\phboxhead{階段（②）}
段の左（$x<0$）では、来る波 $e^{ik_1x}$ とはね返る波 $re^{-ik_1x}$ の重ね合わせ、右（$x>0$）では進む波 $te^{ik_2x}$ だけです（複素数の進む波：原-5 ⑨）。$x=0$ で $\psi$ と $\psi'$ をつなぐと、
\[1+r=t,\qquad ik_1(1-r)=ik_2t\]
で、$t$ を消すと
\[r=\frac{k_1-k_2}{k_1+k_2},\qquad t=\frac{2k_1}{k_1+k_2}\]
です。波-7 ⑤で、ひもの境目で「変位」と「力」をつないで出した $r=(Z_1-Z_2)/(Z_1+Z_2)$ と、まったく同じ計算です。進む波の確率の流れは、速さ（$\propto k$）×$|\psi|^2$ なので、通りぬける割合は
\[\mathcal T=\frac{k_2}{k_1}\,t^2=\frac{4k_1k_2}{(k_1+k_2)^2},\qquad \mathcal R+\mathcal T=1\]
です（波-7 ④の $\dfrac{4Z_1Z_2}{(Z_1+Z_2)^2}$ と同じ形）。

\phboxhead{厚さ $a$ の四角い壁}
壁の左・中・右の3つの所で解を書き、2つの境目で $\psi$ と $\psi'$ をつなぐ（式が4本）と、正確な透過率が出ます。計算は長いので、結果だけを見ます（$\sinh$ は 6-11 ②）。
\pheq{\mathcal T=\frac{1}{1+\dfrac{U_0^{\,2}\sinh^2\kappa a}{4E(U_0-E)}}}{単位なし}
\label{ph:原-6:seikaku}壁が厚い（$\kappa a$ が1より大きい）と、$\sinh\kappa a\approx e^{\kappa a}/2$ なので、分母の1が無視できて、
\[\mathcal T\approx\frac{16E(U_0-E)}{U_0^{\,2}}\,e^{-2\kappa a}\]
になります。前の係数は、$E$ が $U_0$ の半分のとき4で、いつも4以下です。③の $e^{-2\kappa a}$ は、この係数を1とおいた桁の見積もりでした。

\begin{phflow}[境目をつないで、反射と透過へ]
\phflowrow{$\psi$・$\psi'$ がつながる}{段差 $E>U$}{$\mathcal R=\Bigl(\dfrac{k_1-k_2}{k_1+k_2}\Bigr)^2$}
\phflowrow{$\mathcal T=\Bigl[1+\dfrac{U_0^2\sinh^2\kappa a}{4E(U_0-E)}\Bigr]^{-1}$}{厚い壁\\$\kappa a\gg1$}{$\mathcal T\approx e^{-2\kappa a}$（桁）}
\end{phflow}

\begin{bkfigure}
\begin{tikzpicture}
\begin{axis}[bkplot, axis lines=left, ymode=log, width=104mm, height=50mm, xmin=0, xmax=1.25, ymin=1e-6, ymax=2, xtick={0.25,0.5,0.75,1}, ytick={1e-6,1e-4,1e-2,1}, yticklabels={$10^{-6}$,$10^{-4}$,$10^{-2}$,1}, xlabel={厚さ $a$（nm）}, ylabel={$\mathcal T$}]
  \addplot[bkcurve, domain=0:1.2, samples=100] {1/cosh(5.1232*x)^2};
  \addplot[bkcurve2, dash pattern=on 4pt off 1.5pt, domain=0:1.2, samples=100] {exp(-2*5.1232*x)};
  \fill[bkRed] (axis cs:0.5,0.02355) circle (1.8pt) (axis cs:1,1.419e-4) circle (1.8pt);
  \node[txg lab, anchor=west, inner sep=1pt] at (axis cs:0.62,0.05) {正確な式（実線）};
  \node[txg lab, anchor=east, inner sep=1pt] at (axis cs:0.55,2e-4) {$e^{-2\kappa a}$（破線）};
\end{axis}
\end{tikzpicture}
\bkfigcaption{高さ \qty{2}{eV}の壁を \qty{1}{eV}の電子が通りぬける割合（たて軸は10倍ごとの目もり）。どちらも、厚さに対してまっすぐ下がる（指数関数）。正確な値は $e^{-2\kappa a}$ の約4倍。●は例題3。}\label{ph:原-6:fig-toka}
\end{bkfigure}

\phboxhead{波束で見る}
進む波 $e^{ikx}$ は、空間全体に広がった波でした。本物の電子は、ある所に集まった波のかたまり（\textbf{波束}：原-3 ⑦）です。時間をふくむシュレディンガー方程式（原-5 ⑩）を、差分でこまかく時間を進めて解くと、波束が壁に当たるようすを計算できます（図\ref{ph:原-6:fig-hasoku}）。平均のエネルギー \qty{1.5}{eV}の電子の波束を、高さ \qty{2}{eV}・厚さ \qty{0.3}{nm}の壁に当てると、30 fs 後には、波束は2つに分かれ、約30\,\% が壁の向こうへ、約70\,\% がはね返っていました。この割合は、正確な式の $\mathcal T\approx0.30$ と合います。電子1個が「30\,\% だけ」通るのではありません。電子を見つけると、10回に約3回は壁の右で、約7回は左で、まるごと1個見つかるのです。

\begin{bkfigure}
\begin{tikzpicture}
\begin{axis}[bkplot, axis lines=left, width=120mm, height=46mm, xmin=-22, xmax=15, ymin=0, ymax=0.32, xtick={-20,-10,0,10}, ytick={0.1,0.2,0.3}, xlabel={$x$（nm）}, ylabel={$|\psi|^2$（1/nm）}]
  \fill[bkGray!40] (axis cs:0,0) rectangle (axis cs:0.3,0.3);
  \draw[txg psi] plot coordinates {(-22,0) (-21.75,0) (-21.5,0) (-21.25,0) (-21,0) (-20.75,0) (-20.5,0) (-20.25,0) (-20,0) (-19.75,0) (-19.5,0) (-19.25,0) (-19,0) (-18.75,0) (-18.5,0) (-18.25,0) (-18,0.0001) (-17.75,0.0002) (-17.5,0.0003) (-17.25,0.0006) (-17,0.001) (-16.75,0.0018) (-16.5,0.003) (-16.25,0.0048) (-16,0.0076) (-15.75,0.0117) (-15.5,0.0175) (-15.25,0.0255) (-15,0.0361) (-14.75,0.0497) (-14.5,0.0665) (-14.25,0.0866) (-14,0.1096) (-13.75,0.1349) (-13.5,0.1616) (-13.25,0.1882) (-13,0.2132) (-12.75,0.2349) (-12.5,0.2518) (-12.25,0.2624) (-12,0.266) (-11.75,0.2622) (-11.5,0.2514) (-11.25,0.2345) (-11,0.2127) (-10.75,0.1876) (-10.5,0.161) (-10.25,0.1343) (-10,0.109) (-9.75,0.0861) (-9.5,0.0661) (-9.25,0.0493) (-9,0.0358) (-8.75,0.0253) (-8.5,0.0174) (-8.25,0.0116) (-8,0.0075) (-7.75,0.0048) (-7.5,0.0029) (-7.25,0.0018) (-7,0.001) (-6.75,0.0006) (-6.5,0.0003) (-6.25,0.0002) (-6,0.0001) (-5.75,0) (-5.5,0) (-5.25,0) (-5,0) (-4.75,0) (-4.5,0) (-4.25,0) (-4,0) (-3.75,0) (-3.5,0) (-3.25,0) (-3,0) (-2.75,0) (-2.5,0) (-2.25,0) (-2,0) (-1.75,0) (-1.5,0) (-1.25,0) (-1,0) (-0.75,0) (-0.5,0) (-0.25,0) (0.01,0) (0.26,0) (0.51,0) (0.76,0) (1.01,0) (1.26,0) (1.51,0) (1.76,0) (2.01,0) (2.26,0) (2.51,0) (2.76,0) (3.01,0) (3.26,0) (3.51,0) (3.76,0) (4.01,0) (4.26,0) (4.51,0) (4.76,0) (5.01,0) (5.26,0) (5.51,0) (5.76,0) (6.01,0) (6.26,0) (6.51,0) (6.76,0) (7.01,0) (7.26,0) (7.51,0) (7.76,0) (8.01,0) (8.26,0) (8.51,0) (8.76,0) (9.01,0) (9.26,0) (9.51,0) (9.76,0) (10.01,0) (10.26,0) (10.51,0) (10.76,0) (11.01,0) (11.26,0) (11.51,0) (11.76,0) (12.01,0) (12.26,0) (12.51,0) (12.76,0) (13.01,0) (13.26,0) (13.51,0) (13.76,0)};
  \draw[txg psi2] plot coordinates {(-22,0) (-21.75,0) (-21.5,0) (-21.25,0) (-21,0) (-20.75,0) (-20.5,0) (-20.25,0) (-20,0) (-19.75,0) (-19.5,0) (-19.25,0) (-19,0) (-18.75,0) (-18.5,0) (-18.25,0) (-18,0) (-17.75,0) (-17.5,0) (-17.25,0) (-17,0) (-16.75,0.0001) (-16.5,0.0001) (-16.25,0.0001) (-16,0.0002) (-15.75,0.0004) (-15.5,0.0006) (-15.25,0.0009) (-15,0.0013) (-14.75,0.002) (-14.5,0.0029) (-14.25,0.0042) (-14,0.006) (-13.75,0.0083) (-13.5,0.0113) (-13.25,0.0151) (-13,0.0199) (-12.75,0.0258) (-12.5,0.0327) (-12.25,0.0409) (-12,0.0501) (-11.75,0.0604) (-11.5,0.0716) (-11.25,0.0832) (-11,0.0951) (-10.75,0.1068) (-10.5,0.1179) (-10.25,0.1277) (-10,0.136) (-9.75,0.1423) (-9.5,0.1462) (-9.25,0.1476) (-9,0.1464) (-8.75,0.1427) (-8.5,0.1366) (-8.25,0.1285) (-8,0.1187) (-7.75,0.1078) (-7.5,0.0961) (-7.25,0.0842) (-7,0.0725) (-6.75,0.0613) (-6.5,0.0509) (-6.25,0.0416) (-6,0.0333) (-5.75,0.0263) (-5.5,0.0203) (-5.25,0.0154) (-5,0.0115) (-4.75,0.0085) (-4.5,0.0061) (-4.25,0.0043) (-4,0.003) (-3.75,0.0021) (-3.5,0.0014) (-3.25,0.0009) (-3,0.0006) (-2.75,0.0004) (-2.5,0.0002) (-2.25,0.0001) (-2,0.0001) (-1.75,0.0001) (-1.5,0) (-1.25,0) (-1,0) (-0.75,0) (-0.5,0) (-0.25,0) (0.01,0) (0.26,0) (0.51,0) (0.76,0) (1.01,0) (1.26,0) (1.51,0) (1.76,0) (2.01,0) (2.26,0) (2.51,0) (2.76,0.0001) (3.01,0.0001) (3.26,0.0002) (3.51,0.0002) (3.76,0.0004) (4.01,0.0006) (4.26,0.0008) (4.51,0.0012) (4.76,0.0018) (5.01,0.0025) (5.26,0.0035) (5.51,0.0048) (5.76,0.0065) (6.01,0.0085) (6.26,0.0111) (6.51,0.0141) (6.76,0.0176) (7.01,0.0217) (7.26,0.0261) (7.51,0.031) (7.76,0.0361) (8.01,0.0413) (8.26,0.0465) (8.51,0.0514) (8.76,0.0558) (9.01,0.0595) (9.26,0.0623) (9.51,0.0641) (9.76,0.0648) (10.01,0.0644) (10.26,0.0628) (10.51,0.0602) (10.76,0.0567) (11.01,0.0524) (11.26,0.0477) (11.51,0.0426) (11.76,0.0373) (12.01,0.0322) (12.26,0.0272) (12.51,0.0227) (12.76,0.0185) (13.01,0.0149) (13.26,0.0117) (13.51,0.0091) (13.76,0.0069)};
  \node[txg lab, anchor=east] at (axis cs:-13.2,0.22) {はじめ（実線）};
  \node[txg lab, anchor=west, align=left] at (axis cs:-7.2,0.19) {30 fs 後：\\はね返った 70\,\%};
  \node[txg lab, anchor=south] at (axis cs:9.8,0.072) {ぬけた 30\,\%};
  \node[txg lab, anchor=south] at (axis cs:0.15,0.3) {壁};
  \draw[txg arr] (axis cs:-10.5,0.285) -- (axis cs:-5.5,0.285);
\end{axis}
\end{tikzpicture}
\bkfigcaption{電子の波束が壁（灰色）に当たる計算。実線ははじめ、破線は30 fs 後。波束は、はね返る部分と通りぬける部分に分かれる。（時間をふくむ方程式を、差分で時間を進めて解いた。）}\label{ph:原-6:fig-hasoku}
\end{bkfigure}

%% ---------------------------------------------------------------
\bksubsection{8}{調和振動子：$E_n=(n+\tfrac12)\hbar\omega$}
\label{ph:原-6:sec8}
ばねにつけたおもりの位置エネルギーは $U=\tfrac12k_{\mathrm s}x^2$ です。角振動数 $\omega=\sqrt{k_{\mathrm s}/m}$（6-6）を使うと $U=\tfrac12m\omega^2x^2$ と書けます。この位置エネルギーの中の粒を\textbf{調和振動子}といいます。\bkidx{ちょうわしんどうし}{調和振動子}分子の中の原子のゆれ（⑨）や、固体の中の原子のゆれ（熱-15 ⑧）は、谷の底の近くでは、どれもこの形です（力-8：谷の底の近くなら、力はずれに比例する）。

\phboxhead{結果：等しい間かくの準位}
シュレディンガー方程式 $-\dfrac{\hbar^2}{2m}\psi''+\dfrac12m\omega^2x^2\psi=E\psi$ の、遠くで0になる解のエネルギーは、
\pheq{E_n=\Bigl(n+\frac12\Bigr)\hbar\omega\qquad(n=0,\,1,\,2,\,\ldots)}{J ＝ J·s × 1/s}
\label{ph:原-6:chowa}です。すべての解を求める計算は専門の本にゆずり、ここでは、いちばん低い状態を、代入して確かめます。

\phboxhead{いちばん低い状態を代入で確かめる}
$x_0=\sqrt{\hbar/(m\omega)}$ とおき、$\psi_0=e^{-x^2/(2x_0^{\,2})}$ を考えます（ガウスの山の形。9-4）。2回微分すると（合成関数と積の微分、5-6・5-14）、
\[\psi_0'=-\frac{x}{x_0^{\,2}}\,\psi_0,\qquad \psi_0''=\Bigl(\frac{x^2}{x_0^{\,4}}-\frac{1}{x_0^{\,2}}\Bigr)\psi_0\]
です。左辺に入れると、
\[-\frac{\hbar^2}{2m}\Bigl(\frac{x^2}{x_0^{\,4}}-\frac{1}{x_0^{\,2}}\Bigr)\psi_0+\frac12m\omega^2x^2\psi_0\]
です。$x_0^{\,4}=\hbar^2/(m^2\omega^2)$ なので、$\dfrac{\hbar^2}{2mx_0^{\,4}}=\dfrac12m\omega^2$ で、$x^2$ の項は打ち消し合います。残るのは
\[\frac{\hbar^2}{2mx_0^{\,2}}\,\psi_0=\frac{\hbar^2}{2m}\cdot\frac{m\omega}{\hbar}\,\psi_0=\frac12\hbar\omega\,\psi_0\]
で、確かに $E_0=\tfrac12\hbar\omega$ の解です。

\begin{phflow}[方程式から、エネルギーの粒へ]
\phflowrow{$-\dfrac{\hbar^2}{2m}\psi''+U\psi=E\psi$}{$U=\tfrac12m\omega^2x^2$\\遠くで $\psi\to0$}{$E_n=\Bigl(n+\dfrac12\Bigr)\hbar\omega$}
\phflowrow{$E_n=\Bigl(n+\dfrac12\Bigr)\hbar\omega$}{となりの準位の差}{$\Delta\varepsilon=\hbar\omega$（熱-15）}
\end{phflow}

\begin{bkfigure}
\begin{tikzpicture}
\begin{axis}[bkplot, axis lines=left, width=100mm, height=58mm, xmin=-3.4, xmax=3.9, ymin=0, ymax=3.6, xtick={-1,1}, xticklabels={$-x_0$,$x_0$}, ytick={0.5,1.5,2.5}, yticklabels={$\frac12\hbar\omega$,$\frac32\hbar\omega$,$\frac52\hbar\omega$}, xlabel={$x$}, ylabel={$E$}]
  \addplot[bkInk, line width=1.0pt, domain=-2.68:2.68, samples=80] {0.5*x^2};
  \draw[txg lev] (axis cs:-1,0.5) -- (axis cs:1,0.5);
  \draw[txg lev] (axis cs:-1.732,1.5) -- (axis cs:1.732,1.5);
  \draw[txg lev] (axis cs:-2.236,2.5) -- (axis cs:2.236,2.5);
  \addplot[txg psi, domain=-3.2:3.2, samples=120] {0.5+0.42*exp(-x^2/2)};
  \addplot[txg psi, domain=-3.2:3.2, samples=120] {1.5+0.42*1.2*x*exp(-x^2/2)};
  \addplot[txg psi, domain=-3.2:3.2, samples=120] {2.5+0.42*0.7*(2*x^2-1)*exp(-x^2/2)};
  \node[txg lab, anchor=west] at (axis cs:3.3,0.5) {$n=0$};
  \node[txg lab, anchor=west] at (axis cs:3.3,1.5) {$n=1$};
  \node[txg lab, anchor=west] at (axis cs:3.3,2.5) {$n=2$};
  \node[txg lab, anchor=south] at (axis cs:0,2.98) {$U=\tfrac12m\omega^2x^2$};
\end{axis}
\end{tikzpicture}
\bkfigcaption{調和振動子の準位（破線）と波動関数（実線。破線を0として描いた）。準位は $\hbar\omega$ ずつ等しい間かく。破線の両はしは、ふつうの粒が行ける端。波動関数はその外へもしみ出している。}\label{ph:原-6:fig-chowa}
\end{bkfigure}

\begin{phyomi}[$E_n=(n+\tfrac12)\hbar\omega$]
「ゆれのエネルギーは $\hbar\omega$ を1粒として、粒の数 $n$ だけふえる。粒が0個でも $\tfrac12\hbar\omega$ が残る」と読みます。熱-15 ⑧で、ゆれる原子のエネルギーを「粒 $\Delta\varepsilon$ の何個分か」で数え、「粒の大きさは $\omega$ に比例（係数は原-6 で）」と予告しました。その係数が $\hbar$ で、$\Delta\varepsilon=\hbar\omega$ です。
\end{phyomi}

\phboxhead{零点エネルギーと、しみ出し}
いちばん低いエネルギー $\tfrac12\hbar\omega$ を\textbf{零点エネルギー}といいます。\bkidx{れいてんえねるぎー}{零点エネルギー}どんなに冷やしても、ゆれは止まりません。原-5 の箱と同じく、閉じこめられた粒は、不確定性関係（原-3 ⑦）のため、静止していられないからです。熱-15 の最後で「いちばん低いエネルギーは0ではなく $\tfrac12\Delta\varepsilon$。でも比熱は変わらない」と書いたのは、このことです（全部の準位が同じだけ上にずれるだけなので、熱のやりとりには効きません）。

\noindent ふつうの粒がエネルギー $\tfrac12\hbar\omega$ で行ける端は、$\tfrac12m\omega^2x^2=\tfrac12\hbar\omega$ から $x=\pm x_0$ です。ところが、$\psi_0$ はその外にもしみ出していて、$|x|>x_0$ に見つかる確率は約16\,\% もあります。調和振動子の中でも、トンネル効果と同じ「しみこみ」が起きているのです。

\begin{seikaku}
$n\geqq1$ の波動関数は、$e^{-x^2/(2x_0^{\,2})}$ に、エルミート多項式とよばれる多項式をかけた形です（$n=1$ は $x$、$n=2$ は $2x^2/x_0^{\,2}-1$ の形）。全部の解を求める方法と、エルミート多項式の性質は、専門の本で学びます。$n=1$ は章末問題で確かめます。
\end{seikaku}

%% ---------------------------------------------------------------
\bksubsection{9}{分子の振動と赤外分光}
\label{ph:原-6:sec9}
一酸化炭素（CO）のような2つの原子の分子は、2つのおもりを1本のばねでつないだ物とみなせます（波-3 の多自由度振動。原子のばねは力-8）。2つのおもりの質量を $m_1$・$m_2$ とすると、のびちぢみのゆれは、質量が
\pheq{\mu=\frac{m_1m_2}{m_1+m_2}}{kg}
の1つのおもりのゆれと同じになります。$\mu$ を\textbf{換算質量}といい（原-7 ⑪の分子の回転でも使います）、ゆれの角振動数は $\omega=\sqrt{k_{\mathrm s}/\mu}$ です。\bkidx{かんさんしつりょう}{換算質量}

\noindent 分子は、エネルギー $\hbar\omega$ の光子を吸収して、$n=0$ から $n=1$ へ上がります。分子の $\hbar\omega$ は 0.1〜0.5 eV ほどで、波長は約 \qtyrange{2.5}{12}{\micro m}の\textbf{赤外線}です。どの波長の赤外線を吸収するかは、原子の重さとばねの強さで決まるので、分子ごとに違います。赤外線の吸収を調べて物質を見分ける方法を\textbf{赤外分光}といいます。\bkidx{せきがいぶんこう}{赤外分光}『波動』（波-3）で「分子の振動数の測定は、物質を調べる道具」と予告したのは、このことです。

\phboxhead{分光の波数 cm$^{-1}$}
赤外分光では、光の色を「\qty{1}{cm}の中に波がいくつ入るか」で表し、これを（分光の）\textbf{波数} $\tilde\nu$（ニュー・チルダ）とよびます。単位は cm$^{-1}$ です。波長 $\lambda$ と振動数 $f$ とは
\pheq{\tilde\nu=\frac{1}{\lambda},\qquad f=c\,\tilde\nu}{1/cm、Hz ＝ cm/s × 1/cm}
の関係があります（$c=\qty{3.00e10}{cm/s}$ を使う）。CO が吸収する赤外線は $\tilde\nu=\qty{2143}{cm^{-1}}$ で、$\lambda=1/2143\ \text{cm}\approx\qty{4.67}{\micro m}$、$f\approx\qty{6.4e13}{Hz}$ です（例題6）。

\begin{phmijika}[二酸化炭素のセンサーと、温室効果]
部屋の二酸化炭素の濃さをはかるセンサーの多くは、二酸化炭素の分子のゆれが吸収する波長 約 \qty{4.3}{\micro m}の赤外線が、どれだけ弱まるかをはかっています。地球の温室効果も、地面から出る赤外線を、二酸化炭素や水蒸気の分子のゆれが吸収することで起きます。
\end{phmijika}

\begin{motto}[5]
\phboxhead{本物の分子は、少しだけ調和振動子からずれる}
原子のばねは、のばしすぎると弱くなり、最後には切れて（分子がこわれて）しまいます。ちぢめるほうは、原子どうしがぶつかって急に固くなります。だから本物の位置エネルギーは、放物線より右側がゆるやかに開いた形で、準位の間かくは、上へ行くほど少しずつせまくなります（\textbf{非調和}）。分光の測定では、CO の準位の差は、$n=0\to1$ が \qty{2143.2}{cm^{-1}}、$n=1\to2$ が \qty{2116.7}{cm^{-1}}、$n=2\to3$ が \qty{2090.1}{cm^{-1}}で、約 \qty{27}{cm^{-1}}ずつ小さくなります（NIST の分子の定数から計算）。この章で使った \qty{2143}{cm^{-1}}は、いちばん下の2つの準位の差です。いちばん低い準位の近くでは、ずれは1\,\% ほどなので、調和振動子の計算でよく合います。
\end{motto}

%% ---------------------------------------------------------------
\bksubsection{10}{ふだんの温度で、分子の振動が「こおる」わけ}
\label{ph:原-6:sec10}
熱-14 ④で、2つの原子の気体（空気など）の定積モル比熱は、室温で約 $\tfrac52R$ でした。進む動き（3つ）と回る動き（2つ）に $\tfrac12R$ ずつ配られた値で、分子の振動の分が入っていません。「振動は、なぜかエネルギーを受けとらない」というなぞを、ここで解きます。

\noindent CO の振動の $\hbar\omega=\qty{0.266}{eV}$ は、室温（\qty{300}{K}）の $k_{\mathrm B}T=\qty{0.0259}{eV}$ の約10倍です。熱でエネルギーをやりとりするとき、分子が $n=1$ の準位にいる割合は、ボルツマン因子（熱-15 ⑥）で
\[e^{-\hbar\omega/(k_{\mathrm B}T)}=e^{-\theta_{\mathrm E}/T}\approx e^{-10.3}\approx\num{3e-5}\]
しかありません（$\theta_{\mathrm E}=\hbar\omega/k_{\mathrm B}\approx\qty{3080}{K}$）。ほとんど全部の分子が、いちばん下の準位に「こおりついて」います。熱-15 ⑧のアインシュタインの固体と同じ計算で、振動1つ分の比熱は、1 mol あたり
\pheq{C_{\text{振}}=R\,\frac{x^2e^x}{(e^x-1)^2},\qquad x=\frac{\theta_{\mathrm E}}{T}}{J/(mol·K)}
です（熱-15 ⑧の式の、3つの向きの3を1にしたもの）。CO の室温では $C_{\text{振}}\approx0.0036R$ で、ほぼ0です（例題8）。温度が $\theta_{\mathrm E}$ に近づくと、振動にもエネルギーが配られはじめ、高温で $R$ に近づきます（図\ref{ph:原-6:fig-hinetsu}）。\bkidx{しんどうのひねつ}{振動の比熱}

\begin{bkfigure}
\begin{tikzpicture}
\begin{axis}[bkplot, width=104mm, height=44mm, xmin=0, xmax=5300, ymin=0, ymax=1.15, xtick={300,1000,2000,3084,5000}, xticklabels={300,1000,2000,$\theta_{\mathrm E}$,5000}, ytick={0.5,1}, xlabel={$T$（K）}, ylabel={$C_{\text{振}}/R$}]
  \addplot[bkcurve, domain=120:5200, samples=150] {(3084/x)^2*exp(3084/x)/(exp(3084/x)-1)^2};
  \fill[bkRed] (axis cs:300,0.0036) circle (1.8pt) (axis cs:1000,0.478) circle (1.8pt);
\end{axis}
\end{tikzpicture}
\bkfigcaption{CO の振動の比熱（1 mol あたり、$R$ を単位に）。●は \qty{300}{K}（例題8）と \qty{1000}{K}（じぶんでとこう8）。室温ではほぼ0で、振動は「こおっている」。}\label{ph:原-6:fig-hinetsu}
\end{bkfigure}

\noindent 回る動きがこおらないのは、回転のエネルギーの段（CO で約 \qty{0.0005}{eV}：原-7 ⑪）が $k_{\mathrm B}T$ よりずっと小さいからです。とびとびのエネルギーの段の大きさと $k_{\mathrm B}T$ を比べれば、どの動きにエネルギーが配られるかがわかります。19世紀の物理学者を悩ませた比熱のなぞ（熱-14）は、量子力学で解けたのです。

%% ---------------------------------------------------------------
\bksubsection{11}{2つの井戸とトンネル分裂：アンモニア}
\label{ph:原-6:sec11}
同じ形の井戸を2つ、うすい壁をはさんで並べます。壁が厚ければ、左の井戸の状態 $\psi_{\text{左}}$ と右の井戸の状態 $\psi_{\text{右}}$ は、同じエネルギー $E_0$ をもちます。壁がうすいと、トンネル効果で行き来できるようになり、2つの状態が混ざります。左と右の状態を矢印 $(1,0)$・$(0,1)$ のように表し、行き来の強さを $A$（$>0$）とすると、エネルギーの固有値問題は $2\times2$ の行列
\[\begin{pmatrix}E_0&-A\\-A&E_0\end{pmatrix}\]
の固有値問題になります（7-8）。固有値は $E_0-A$ と $E_0+A$、固有ベクトルは $(1,1)/\sqrt2$（左右の和：対称）と $(1,-1)/\sqrt2$（左右の差：反対称）です。1つだった準位が、差 $2A$ の2つに分かれます。これを\textbf{トンネル分裂}といいます。\bkidx{とんねるぶんれつ}{トンネル分裂}

\begin{bkfigure}
\begin{tikzpicture}
\begin{axis}[bkplot, axis x line=middle, axis y line=none, width=104mm, height=46mm, xmin=-2.7, xmax=2.7, ymin=-1.0, ymax=1.15, xtick={-1,1}, ytick=\empty, xlabel={$x$（nm）}]
  \fill[bkGray!25] (axis cs:-2.6,-0.98) rectangle (axis cs:-1.25,1.1) (axis cs:-0.25,-0.98) rectangle (axis cs:0.25,1.1) (axis cs:1.25,-0.98) rectangle (axis cs:2.6,1.1);
  \draw[txg psi] plot coordinates {(-2.6,0.001) (-2.54,0.001) (-2.48,0.001) (-2.42,0.002) (-2.36,0.002) (-2.3,0.003) (-2.24,0.004) (-2.18,0.005) (-2.12,0.006) (-2.06,0.008) (-2,0.011) (-1.94,0.015) (-1.88,0.019) (-1.82,0.025) (-1.76,0.033) (-1.7,0.044) (-1.64,0.058) (-1.58,0.077) (-1.52,0.101) (-1.46,0.134) (-1.4,0.177) (-1.34,0.234) (-1.28,0.309) (-1.22,0.402) (-1.16,0.493) (-1.1,0.574) (-1.04,0.646) (-0.98,0.707) (-0.92,0.756) (-0.86,0.792) (-0.8,0.814) (-0.74,0.823) (-0.68,0.817) (-0.62,0.798) (-0.56,0.764) (-0.5,0.718) (-0.44,0.66) (-0.38,0.59) (-0.32,0.51) (-0.26,0.421) (-0.2,0.339) (-0.14,0.283) (-0.08,0.249) (-0.02,0.235) (0.04,0.239) (0.1,0.261) (0.16,0.304) (0.22,0.37) (0.28,0.459) (0.34,0.544) (0.4,0.62) (0.46,0.685) (0.52,0.739) (0.58,0.78) (0.64,0.807) (0.7,0.821) (0.76,0.821) (0.82,0.807) (0.88,0.779) (0.94,0.737) (1,0.683) (1.06,0.618) (1.12,0.541) (1.18,0.456) (1.24,0.363) (1.3,0.275) (1.36,0.208) (1.42,0.158) (1.48,0.119) (1.54,0.09) (1.6,0.068) (1.66,0.052) (1.72,0.039) (1.78,0.03) (1.84,0.023) (1.9,0.017) (1.96,0.013) (2.02,0.01) (2.08,0.007) (2.14,0.006) (2.2,0.004) (2.26,0.003) (2.32,0.002) (2.38,0.002) (2.44,0.001) (2.5,0.001) (2.56,0.001)};
  \draw[txg psi2] plot coordinates {(-2.6,-0.001) (-2.54,-0.001) (-2.48,-0.001) (-2.42,-0.002) (-2.36,-0.002) (-2.3,-0.003) (-2.24,-0.004) (-2.18,-0.005) (-2.12,-0.007) (-2.06,-0.009) (-2,-0.012) (-1.94,-0.016) (-1.88,-0.021) (-1.82,-0.028) (-1.76,-0.037) (-1.7,-0.049) (-1.64,-0.064) (-1.58,-0.085) (-1.52,-0.111) (-1.46,-0.147) (-1.4,-0.193) (-1.34,-0.254) (-1.28,-0.334) (-1.22,-0.435) (-1.16,-0.53) (-1.1,-0.616) (-1.04,-0.69) (-0.98,-0.751) (-0.92,-0.798) (-0.86,-0.83) (-0.8,-0.846) (-0.74,-0.846) (-0.68,-0.83) (-0.62,-0.799) (-0.56,-0.753) (-0.5,-0.692) (-0.44,-0.618) (-0.38,-0.533) (-0.32,-0.438) (-0.26,-0.334) (-0.2,-0.234) (-0.14,-0.153) (-0.08,-0.082) (-0.02,-0.018) (0.04,0.044) (0.1,0.11) (0.16,0.185) (0.22,0.273) (0.28,0.378) (0.34,0.478) (0.4,0.57) (0.46,0.65) (0.52,0.719) (0.58,0.774) (0.64,0.814) (0.7,0.839) (0.76,0.848) (0.82,0.841) (0.88,0.818) (0.94,0.78) (1,0.727) (1.06,0.661) (1.12,0.582) (1.18,0.492) (1.24,0.392) (1.3,0.298) (1.36,0.227) (1.42,0.172) (1.48,0.131) (1.54,0.099) (1.6,0.076) (1.66,0.057) (1.72,0.044) (1.78,0.033) (1.84,0.025) (1.9,0.019) (1.96,0.015) (2.02,0.011) (2.08,0.008) (2.14,0.006) (2.2,0.005) (2.26,0.004) (2.32,0.003) (2.38,0.002) (2.44,0.002) (2.5,0.001) (2.56,0.001)};
  \node[txg lab] at (axis cs:-1.95,0.55) {対称（実線）};
  \node[txg lab] at (axis cs:-1.95,-0.55) {反対称（破線）};
  \node[txg lab, anchor=north] at (axis cs:0,1.1) {壁};
\end{axis}
\end{tikzpicture}
\bkfigcaption{深さ \qty{1}{eV}・幅 \qty{1}{nm}の井戸を2つ、厚さ \qty{0.5}{nm}の壁（灰色の間）をはさんで並べたときの、低い2つの状態（差分法で計算）。エネルギーは \qty{0.181}{eV}と \qty{0.200}{eV}で、差は \qty{19}{meV}。}\label{ph:原-6:fig-futatsu}
\end{bkfigure}

\noindent 差分法（原-5 ⑫）で計算すると、図\ref{ph:原-6:fig-futatsu}の2つの井戸の分裂は \qty{18.7}{meV}、壁を \qty{1}{nm}に厚くすると \qty{1.87}{meV}で、10分の1になります。分裂は、壁をすりぬける割合と同じように、厚さで指数関数的に小さくなるのです。

\noindent はじめ電子を左の井戸に入れると、それは対称と反対称の重ね合わせ $(\psi_{\text{対}}+\psi_{\text{反}})/\sqrt2$ です。原-5 ⑩と同じく、2つの矢印が違う速さで回るので、電子は周期 $h/(2A)$ で左右の井戸を行き来します。

\phboxhead{アンモニアの分子}
アンモニア（NH$_3$）の分子は、3つの水素原子が作る三角形の面の上か下に、窒素原子がいる、三角すいの形です。窒素が面の上にいる形と下にいる形は、2つの井戸の底にあたり、間には面を通りぬける壁があります。窒素はこの壁をトンネル効果で行き来し、準位が2つに分かれます。その差は、振動数 \qty{23.87}{GHz}（波長約 \qty{1.26}{cm}）のマイクロ波の光子のエネルギーにあたります（例題11）。1954年、タウンズらは、上の準位の分子だけを電場でより分けて集め、この振動数のマイクロ波を誘導放出で強めました。これが最初の\textbf{メーザー}で（原-4 ⑫）、レーザーの先がけとなり、1964年のノーベル物理学賞につながりました。

%% ---------------------------------------------------------------
\bksubsection{12}{目に見える大きさのトンネル効果}
\label{ph:原-6:sec12}
トンネル効果は、電子や原子のような小さな粒のものだと考えられてきました。ところが、1984〜1985年、クラーク、デヴォレ、マーティニスは、超伝導の金属で作った、手でつまめるほどの大きさのチップの上の電気の回路で、回路全体が1つの粒のようにふるまい、閉じこめられた状態からトンネル効果でぬけ出すことと、そのエネルギーがとびとびであることを、実験で示しました。たくさんの電子が1つにまとまって動く超伝導（原-8 ⑫）だからこそ起きることです。この研究は、2025年のノーベル物理学賞を受けました。同じしくみの回路は、いま、量子コンピュータの部品（量子ビット）として使われています。

\begin{phmada}[すりぬけるのに、どれだけ時間がかかる？]
トンネル効果で壁をすりぬける電子は、壁の中に「どれだけの時間」いるのでしょう。じつは、何をはかって「時間」とよぶかで、答えが変わります。強いレーザーで原子から電子を引き出す実験（アト秒の時計、2019年、水素の原子）では、おくれは1.8アト秒（$\qty{1.8e-18}{s}$）より小さく、ほぼ0と報告されました。いっぽう、冷やした原子の向き（スピン）が壁の中でだけ回るようにして、回った角度を時計にした実験（2020年）では、0でない時間がはかられました。2つの実験は違う量をはかっているという考えが出されていますが、「トンネルの時間」をどう決めればよいかは、今も議論が続いています。
\end{phmada}

\begin{column}{原子で字を書いた}
走査トンネル顕微鏡（④）の針は、原子を見るだけでなく、動かすこともできます。1990年、アメリカの研究者は、約 \qty{4}{K}に冷やしたニッケルの表面の上で、針を使ってキセノンの原子を1個ずつ動かし、35個の原子で会社の名前の3文字を書きました。人間が原子を1個ずつ並べて形を作った、最初の例の1つです。いまでは、原子を並べて小さな「囲い」を作り、その中に閉じこめられた電子の波（原-5 の箱の電子と同じ定在波）を、STM で写真にとることもできます。
\end{column}

\begin{senmonsho}[半導体工学の専門書の書き方]
\begin{bkhonbun}
WKB近似によれば，透過係数は\[T\simeq\exp\Bigl(-2\int_{x_1}^{x_2}\kappa(x)\,dx\Bigr),\qquad \kappa(x)=\frac{\sqrt{2m^*(U(x)-E)}}{\hbar}\]で与えられる．三角形障壁に対してはファウラー・ノルドハイム電流\[J\propto\mathcal E_{\mathrm{ox}}^{\,2}\exp\Bigl(-\frac{B}{\mathcal E_{\mathrm{ox}}}\Bigr)\]を得る．
\end{bkhonbun}
\bkyaku{壁の高さが場所で変わるときは、$e^{-2\kappa a}$ の $\kappa a$ を、場所ごとの $\kappa$ の積分（たし合わせ）にかえる（WKB 近似）。$m^*$ は有効質量（原-5 ⑤）。ななめ（三角形）の壁では、電流が酸化膜の電場 $\mathcal E_{\mathrm{ox}}$ で急に増える式になる（④の書きこみ）。$\simeq$ は「ほぼ等しい」。}
\end{senmonsho}

\begin{phsaki}{トンネル効果と分子のゆれの先へ}
専門の本では、壁の形が場所で変わるときの近似（WKB 近似）、共鳴して通りぬけやすくなる2重の壁（共鳴トンネルダイオード）、電子のスピンで通りぬけやすさが変わる磁気トンネル接合（ハードディスクの読み取りや MRAM）を学びます。分子のゆれは、調和振動子からのずれ（非調和）や、回転と組み合わさったスペクトルとして、分子分光学で学びます。超伝導の回路のトンネル効果は、量子コンピュータの工学につながります。
\end{phsaki}

\begin{taisetsu}
\def\tsl#1{\makebox[8.4\zw][l]{\sffamily\bfseries\small\color{\bkTT}#1}}
\tsl{壁の中}$\psi\propto e^{-\kappa x}$、$\kappa=\dfrac{\sqrt{2m(U_0-E)}}{\hbar}\approx5.12\sqrt{U_0-E}$〔nm$^{-1}$〕\\[3pt]
\tsl{すりぬける割合}$\mathcal T\approx e^{-2\kappa a}$（厚さ2倍で2乗）\\[3pt]
\tsl{階段}$\mathcal R=\Bigl(\dfrac{k_1-k_2}{k_1+k_2}\Bigr)^2$（波-7 の $Z$ を $k$ に）\\[3pt]
\tsl{有限の井戸}偶 $z\tan z=\sqrt{z_0^{\,2}-z^2}$、奇 $-z/\tan z=\sqrt{z_0^{\,2}-z^2}$\\[3pt]
\tsl{調和振動子}$E_n=\Bigl(n+\dfrac12\Bigr)\hbar\omega$、$\psi_0=e^{-x^2/(2x_0^2)}$、$x_0=\sqrt{\dfrac{\hbar}{m\omega}}$\\[3pt]
\tsl{分子の振動}$\omega=\sqrt{k_{\mathrm s}/\mu}$、$\mu=\dfrac{m_1m_2}{m_1+m_2}$、$f=c\tilde\nu$
\end{taisetsu}

\begin{kotoba}
\small
\bkword{トンネル効果}{エネルギーが足りなくても、うすい壁を波としてすりぬけること。}
\bkword{透過率・反射率}{壁や段を通りぬける割合・はね返る割合。$\mathcal T$・$\mathcal R$。}
\bkword{浮遊ゲート}{フラッシュメモリで電子を閉じこめる、絶縁膜で囲まれた導体。}
\bkword{走査トンネル顕微鏡}{トンネル電流で、表面の原子を見る顕微鏡（STM）。}
\bkword{調和振動子}{$U=\tfrac12k_{\mathrm s}x^2$ の中の粒。準位は等しい間かく $\hbar\omega$。}
\bkword{零点エネルギー}{いちばん低い状態にも残るエネルギー。調和振動子では $\tfrac12\hbar\omega$。}
\bkword{換算質量}{2つの物の相対的なゆれを、1つのおもりにおきかえたときの質量。}
\bkword{赤外分光}{分子のゆれが吸収する赤外線の波長で、物質を調べる方法。}
\bkword{トンネル分裂}{2つの井戸の間の行き来で、準位が2つに分かれること。}
\end{kotoba}

\bkkigocards{\bkkigocard{$\kappa$}{カッパ}{減衰の定数〔1/m〕。壁の中の波が $e^{-\kappa x}$ で弱まる速さ。$\sqrt{2m(U_0-E)}/\hbar$。原-8 の熱伝導率 $\kappa$ とは別}{decay constant}{原-6}}
{\bkkigocard{$\mathcal T$・$\mathcal R$}{ティー・アール（筆記体）}{透過率・反射率（単位なし）。温度 $T$・気体定数 $R$ と分けるため筆記体で書く}{transmission, reflection}{原-6}}
\bkkigocards{\bkkigocard{$\mu$}{ミュー}{換算質量〔kg〕。$m_1m_2/(m_1+m_2)$（原-7 ⑪でも）。減弱係数（原-13）・透磁率 $\mu_0$ とは別}{reduced mass}{原-6}}
{\bkkigocard{$\tilde\nu$}{ニュー・チルダ}{分光の波数〔cm$^{-1}$〕。1 cm に入る波の数。$f=c\tilde\nu$}{wavenumber}{原-6}}

\begin{column}{電子レンジは、分子をゆらしているの？}
電子レンジの電磁波（周波数 \qty{2.45}{GHz}）の光子1個のエネルギー $hf$ は、約 \qty{1.0e-5}{eV}です。CO の振動の $\hbar\omega$（\qty{0.27}{eV}）の約2万6千分の1しかなく、振動の準位を上げることはできません。水の分子は電気のかたより（＋と－の向き）をもっていて、1秒に24億5千万回向きを変える電場に引っぱられて回ろうとしますが、周りの分子とぶつかってうまく回れず、その動きが熱になります（誘電加熱）。分子の振動の吸収で温める赤外線のヒーター（⑨）とは、違うしくみです。
\end{column}

%% ======================== 例題 ========================
\Needspace{10\baselineskip}\bksection{reidai}{例題でたしかめよう}
\label{ph:原-6:ex1}
\begin{phreidai}[USB メモリの電子は、何年もとどまる？]
フラッシュメモリの浮遊ゲートを囲む酸化膜を、高さ（電子のエネルギーとの差）\qty{3.2}{eV}の四角い壁とします。厚さ 2・5・\qty{8}{nm}のとき、すりぬける割合 $e^{-2\kappa a}$ を求めましょう。また、電子が1秒に $10^{13}$ 回壁に当たるとすると、\qty{2}{nm}の壁では、電子がぬけるまでにどれくらいかかりますか。
\phbamen
はてな？：記録の電子が、壁にどれだけしっかり閉じこめられているかを見積もる場面です。
\phmushi
電子の有効質量は自由な電子と同じとし、電場はかかっていない（壁は四角）とします。1秒に当たる回数は、桁の見当として置いた値です。
\phzu
図\ref{ph:原-6:fig-kabe}(ア)。$U_0-E=\qty{3.2}{eV}$、$a=2,\ 5,\ 8$ nm。
\phriyu
壁の中で $\psi\propto e^{-\kappa x}$、確率はその2乗なので $e^{-2\kappa a}$（③）。
\phtateru
\[\kappa=5.12\sqrt{3.2},\qquad \mathcal T\approx e^{-2\kappa a},\qquad \text{ぬけるまでの時間}\approx\frac{1}{10^{13}\times\mathcal T}\]
\phtoku
$\kappa=5.12\times1.789\approx\qty{9.17}{nm^{-1}}$。厚さ 2・5・\qty{8}{nm}で、
\[e^{-36.7}\approx\num{1.2e-16},\qquad e^{-91.7}\approx\num{1.6e-40},\qquad e^{-146.6}\approx\num{2.1e-64}\]
\qty{2}{nm}では、1秒にぬける見こみは $10^{13}\times1.2\times10^{-16}=\num{1.2e-3}$ で、ぬけるまで約 $1\div0.0012\approx\qty{830}{s}$（約14分）。
\phtashikame
\qty{8}{nm}なら、$10^{13}\times2\times10^{-64}$ で1秒に $10^{-51}$ 回ほどで、宇宙の年齢（約 $\qty{4e17}{s}$）待ってもぬけません。厚さが4倍で、割合は $10^{-16}$ の4乗ほどになっています（③の2倍にすると2乗）。うすすぎる膜では、記録が数分で消えてしまいます。
\kotae $\num{1.2e-16}$、$\num{1.6e-40}$、$\num{2.1e-64}$。\qty{2}{nm}では約14分でぬける
\end{phreidai}
\begin{jibun}
高さの差が \qty{1}{eV}の壁で、厚さ \qty{1}{nm}と \qty{2}{nm}のとき、すりぬける割合 $e^{-2\kappa a}$ を求めましょう。
\bkjans{\qty{1}{nm}：約 $\num{3.6e-5}$、\qty{2}{nm}：約 $\num{1.3e-9}$}
\end{jibun}

\label{ph:原-6:ex2}
\begin{phreidai}[段差ではね返る電子]
エネルギー \qty{2}{eV}の電子が、高さ \qty{1}{eV}の段差に来ました。はね返る割合（反射率）を求め、波-7 の式と比べましょう。
\phbamen
ふつうの粒なら全部上がれる段で、電子の波の一部がはね返る場面です（②）。
\phmushi
段は急な（はばのない）段差とします。
\phzu
図\ref{ph:原-6:fig-kaidan}。$E=\qty{2}{eV}$、$U=\qty{1}{eV}$。
\phriyu
段の左右で波数が変わり、境目で $\psi$ と $\psi'$ がつながる。結果は②の $\mathcal R$。
\phtateru
\[\frac{k_2}{k_1}=\sqrt{\frac{E-U}{E}},\qquad \mathcal R=\Bigl(\frac{1-k_2/k_1}{1+k_2/k_1}\Bigr)^2\]
\phtoku
$k_2/k_1=\sqrt{0.5}\approx0.707$。
\[\mathcal R=\Bigl(\frac{0.293}{1.707}\Bigr)^2\approx0.0294\]
\phtashikame
約3\,\% がはね返り、約97\,\% が通りぬけます（反射率と透過率の和は1）。波-7 の2本のひもで、インピーダンスの比が $Z_2/Z_1=0.707$ のときと同じ値です。
\kotae 約0.029（約3\,\%）
\end{phreidai}
\begin{jibun}
エネルギー \qty{1.5}{eV}の電子が、高さ \qty{1}{eV}の段差に来たときの反射率を求めましょう。
\bkjans{$k_2/k_1=\sqrt{0.5/1.5}\approx0.577$。$\mathcal R=(0.423/1.577)^2\approx0.072$}
\end{jibun}

\label{ph:原-6:ex3}
\begin{phreidai}[うすい壁の正確な透過率]
高さ \qty{2}{eV}の壁に、エネルギー \qty{1}{eV}の電子が来ます。壁の厚さが \qty{0.5}{nm}と \qty{1.0}{nm}のとき、正確な透過率と、③の近似 $e^{-2\kappa a}$ を求め、比べましょう。
\phbamen
桁の見積もり $e^{-2\kappa a}$ が、どれくらい正しいかを確かめる場面です（⑦）。
\phmushi
壁は四角とします。
\phzu
図\ref{ph:原-6:fig-toka}。$U_0=2$、$E=1$、$U_0-E=1$〔eV〕。
\phriyu
正確な式は⑦。$U_0=2E$ なので、分母の $\sinh^2$ の前の係数は $4/4=1$ になり、$\mathcal T=1/(1+\sinh^2\kappa a)$。
\phtateru
\[\kappa=\qty{5.12}{nm^{-1}},\qquad \mathcal T=\frac{1}{1+\sinh^2\kappa a},\qquad \mathcal T_{\text{近}}=e^{-2\kappa a}\]
\phtoku
\qty{0.5}{nm}：$\sinh2.56\approx6.44$、$\mathcal T=1/(1+6.44^2)\approx0.024$、近似 $e^{-5.12}\approx0.0060$。\\
\qty{1.0}{nm}：$\sinh5.12\approx83.9$、$\mathcal T\approx\num{1.4e-4}$、近似 $e^{-10.24}\approx\num{3.6e-5}$。
\phtashikame
正確な値は近似の約4倍で、⑦の係数 $16E(U_0-E)/U_0^{\,2}=4$ と合います（$4\times0.0060=0.024$）。近似は桁で合い、厚さを2倍にすると割合が約2乗になることも同じです（$0.024^2\approx\num{6e-4}$ の4分の1ほど。前の係数4が1回分しかかからないため）。
\kotae \qty{0.5}{nm}：0.024（近似 0.0060）、\qty{1.0}{nm}：$\num{1.4e-4}$（近似 $\num{3.6e-5}$）
\end{phreidai}
\begin{jibun}
同じ壁と電子で、厚さ \qty{0.3}{nm}のときの正確な透過率と近似 $e^{-2\kappa a}$ を求めましょう（$\sinh1.537\approx2.23$）。
\bkjans{正確：約0.17、近似：約0.046（うすい壁では、4倍の0.19とも少しずれる）}
\end{jibun}

\label{ph:原-6:ex4}
\begin{phreidai}[STM の針を 0.1 nm 近づける]
走査トンネル顕微鏡で、針と表面の間の壁の高さ（仕事関数）を約 \qty{4.5}{eV}とします。針を表面に \qty{0.1}{nm}近づけると、トンネル電流は何倍になりますか。
\phbamen
STM が原子1個の高さの違いを見分けられるわけを、数で確かめる場面です（④）。
\phmushi
電子のエネルギーは壁の高さに比べて小さく（$U_0-E\approx\qty{4.5}{eV}$）、電流はすりぬける割合に比例するとします。
\phzu
図\ref{ph:原-6:fig-stm}。すき間 $d$ が $d-0.1$ nm になる。
\phriyu
電流 $\propto e^{-2\kappa d}$（③）。すき間が $\Delta d$ 小さくなると、$e^{2\kappa\Delta d}$ 倍。
\phtateru
\[\kappa=5.12\sqrt{4.5},\qquad \text{倍率}=e^{2\kappa\times0.1}\]
\phtoku
$\kappa=5.12\times2.121\approx\qty{10.9}{nm^{-1}}$。$2\kappa\times0.1\approx2.17$ で、$e^{2.17}\approx8.8$。
\phtashikame
原子の大きさの半分ほどの \qty{0.1}{nm}で、電流が約9倍も変わります。電流を1\,\% の精度ではかれば、\qty{0.001}{nm}ほどの高さの違いがわかる計算です。
\kotae 約8.8倍（約9倍）
\end{phreidai}
\begin{jibun}
同じ STM で、電流をちょうど10倍にするには、針を何 nm 近づければよいですか（$\ln10\approx2.30$）。
\bkjans{$\Delta d=\ln10/(2\kappa)\approx\qty{0.11}{nm}$}
\end{jibun}

\label{ph:原-6:ex5}
\begin{phreidai}[深さ 1 eV・幅 1 nm の井戸]
深さ \qty{1}{eV}、幅 \qty{1}{nm}の井戸に閉じこめられた電子の準位を、グラフの交点で求めましょう。同じ幅の無限の井戸（原-5 例題1）と比べましょう。
\phbamen
壁の高さに限りがある井戸の、とびとびの準位を求める場面です（⑥）。
\phmushi
井戸は四角とし、電子は1次元を動くとします。
\phzu
図\ref{ph:原-6:fig-kouten}。$U_0=\qty{1}{eV}$、$a=\qty{1}{nm}$。
\phriyu
境目で $\psi$ と $\psi'$ をつなぐと、偶・奇の式になる（⑥）。式の変形では解けないので、グラフの交点（表計算）で求める。
\phtateru
\[z_0=\frac a2\times5.12\sqrt{U_0},\qquad z\tan z=\sqrt{z_0^{\,2}-z^2},\qquad -\frac{z}{\tan z}=\sqrt{z_0^{\,2}-z^2}\]
\phtoku
$z_0=0.5\times5.12\times1\approx2.56$。$z_0$ は $\pi/2=1.57$ より大きく $\pi=3.14$ より小さいので、交点は偶1つと奇1つ。表計算で $z$ を変えてさがすと、偶は $z\approx1.119$、奇は $z\approx2.147$。\\
$E=0.0381\times(2z/a)^2$ より、偶：$0.0381\times2.238^2\approx\qty{0.19}{eV}$。\\
奇：$0.0381\times4.295^2\approx\qty{0.70}{eV}$。
\phtashikame
無限の井戸の \qty{0.376}{eV}・\qty{1.50}{eV}より低い（しみ出して、実際の幅が広がった）。どちらも深さ \qty{1}{eV}より低く、閉じこめられています。
\kotae 閉じこめられた状態は2つ。底から約 \qty{0.19}{eV}と約 \qty{0.70}{eV}
\end{phreidai}
\begin{jibun}
深さ \qty{0.3}{eV}、幅 \qty{1}{nm}の井戸では、閉じこめられた状態はいくつありますか。いちばん低い準位も求めましょう。
\bkjans{$z_0\approx1.40<\pi/2$ なので1つ。約 \qty{0.12}{eV}}
\end{jibun}

\label{ph:原-6:ex6}
\begin{phreidai}[CO の振動のエネルギー]
一酸化炭素（CO）の分子は、波数 \qty{2143}{cm^{-1}}の赤外線を吸収して、振動の $n=0$ から $n=1$ へ上がります。振動数 $f$、$\hbar\omega$、零点エネルギーを求めましょう。
\phbamen
分子の振動を、調和振動子として数で読む場面です（⑧⑨）。
\phmushi
振動は調和振動子とします（大きなゆれでのずれ＝非調和を無視）。
\phzu
$\tilde\nu=\qty{2143}{cm^{-1}}$、$c=\qty{3.00e10}{cm/s}$、$h=\qty{4.14e-15}{eV.s}$。
\phriyu
吸収する光子のエネルギー $hf$ が、準位の差 $\hbar\omega$ に等しい（⑧）。
\phtateru
\[f=c\tilde\nu,\qquad \hbar\omega=hf,\qquad E_0=\frac12\hbar\omega\]
\phtoku
\[f=3.00\times10^{10}\times2143\approx\qty{6.43e13}{Hz}\]
$\hbar\omega=hf=4.14\times10^{-15}\times6.43\times10^{13}\approx\qty{0.266}{eV}$。零点エネルギーは、その半分で約 \qty{0.133}{eV}。
\phtashikame
室温の $k_{\mathrm B}T$（\qty{0.026}{eV}）の約10倍です（⑩で使う）。分子の振動の 0.1〜0.5 eV の範囲に入っています。
\kotae $f\approx\qty{6.4e13}{Hz}$、$\hbar\omega\approx\qty{0.27}{eV}$、零点エネルギー約 \qty{0.13}{eV}
\end{phreidai}
\begin{jibun}
塩化水素（HCl）は、波数 \qty{2885}{cm^{-1}}の赤外線を吸収します。$f$、$\hbar\omega$、零点エネルギーを求めましょう。
\bkjans{$f\approx\qty{8.66e13}{Hz}$、$\hbar\omega\approx\qty{0.36}{eV}$、零点エネルギー約 \qty{0.18}{eV}}
\end{jibun}

\label{ph:原-6:ex7}
\begin{phreidai}[CO の原子のばね定数]
例題6の CO の振動から、炭素と酸素をつなぐ「原子のばね」のばね定数を求めましょう。換算質量は \qty{6.86}{u}（$\qty{1}{u}=\qty{1.66e-27}{kg}$）とします。
\phbamen
物理シリーズの『力学』（力-8）で「原子の間の力はばね」と見たばね定数を、分子の光から求める場面です。
\phmushi
調和振動子とします。分子の回転は考えません。
\phzu
$\mu=6.86\times1.66\times10^{-27}\approx\qty{1.14e-26}{kg}$、$f=\qty{6.43e13}{Hz}$。
\phriyu
2つの原子の振動は、換算質量 $\mu$ の1つのおもりの振動と同じ（⑨）。$\omega=\sqrt{k_{\mathrm s}/\mu}$ を $k_{\mathrm s}$ について解く。
\phtateru
\[k_{\mathrm s}=\mu\omega^2=\mu(2\pi f)^2\]
\phtoku
$2\pi f\approx\qty{4.04e14}{rad/s}$。
\[k_{\mathrm s}=1.14\times10^{-26}\times(4.04\times10^{14})^2\approx\qty{1.9e3}{N/m}\]
\phtashikame
\qty{1}{nm}のびるのに約 $\qty{1.9e-6}{N}$ の力がいる、ということです。ふつうのばね（数十 N/m）よりずっと固いばねが、原子の間にあります。単位は kg/s$^2$＝N/m で合っています。
\kotae 約 $\qty{1.9e3}{N/m}$
\end{phreidai}
\begin{jibun}
HCl（換算質量 \qty{0.980}{u}、$f=\qty{8.65e13}{Hz}$）のばね定数を求めましょう。
\bkjans{$\mu\approx\qty{1.63e-27}{kg}$、$\omega\approx\qty{5.44e14}{rad/s}$ で、$k_{\mathrm s}=\mu\omega^2\approx\qty{4.8e2}{N/m}$}
\end{jibun}

\label{ph:原-6:ex8}
\begin{phreidai}[CO の振動は、室温で比熱に効くか]
CO の振動の $\hbar\omega=\qty{0.266}{eV}$ として、アインシュタイン温度 $\theta_{\mathrm E}$ と、\qty{300}{K}での振動の比熱 $C_{\text{振}}/R$ を求めましょう。
\phbamen
熱-14 の「振動はエネルギーを受けとらない」なぞを、数で確かめる場面です（⑩）。
\phmushi
分子どうしの力と、回転と振動の結びつきは考えません。
\phzu
図\ref{ph:原-6:fig-hinetsu}。$k_{\mathrm B}=\qty{8.62e-5}{eV/K}$。
\phriyu
振動1つの比熱は、熱-15 ⑧の式で $x=\theta_{\mathrm E}/T$ とおいたもの（⑩）。
\phtateru
\[\theta_{\mathrm E}=\frac{\hbar\omega}{k_{\mathrm B}},\qquad \frac{C_{\text{振}}}{R}=\frac{x^2e^x}{(e^x-1)^2}\]
\phtoku
$\theta_{\mathrm E}=0.266\div(8.62\times10^{-5})\approx\qty{3080}{K}$。$x=3080\div300\approx10.3$。$e^{10.3}\approx2.9\times10^4$ なので、
\[\frac{C_{\text{振}}}{R}\approx\frac{10.3^2\times2.9\times10^4}{(2.9\times10^4)^2}\approx\frac{106}{2.9\times10^4}\approx0.0036\]
\phtashikame
\qty{1}{mol}あたり約 $0.0036\times8.31\approx\qty{0.03}{J/K}$ で、$\tfrac52R\approx\qty{20.8}{J/(mol.K)}$ に比べてほぼ0。振動は「こおって」います。
\kotae $\theta_{\mathrm E}\approx\qty{3080}{K}$、$C_{\text{振}}/R\approx0.0036$（ほぼ0）
\end{phreidai}
\begin{jibun}
CO の振動の比熱 $C_{\text{振}}/R$ を、\qty{1000}{K}で求めましょう（$e^{3.08}\approx21.8$）。
\bkjans{$x\approx3.08$。$3.08^2\times21.8\div20.8^2\approx0.48$}
\end{jibun}

\label{ph:原-6:ex9}
\begin{phreidai}[HCl が吸収する赤外線の波長]
塩化水素（HCl）の振動は、波数 \qty{2885}{cm^{-1}}です。吸収する赤外線の波長を求めましょう。
\phbamen
赤外分光で、分子が吸収する光の波長を読む場面です（⑨。波-3 の予告）。
\phmushi
回転による細かい線の分かれは考えません。
\phzu
$\tilde\nu=\qty{2885}{cm^{-1}}$。
\phriyu
波数は1 cm に入る波の数なので、$\lambda=1/\tilde\nu$（⑨）。
\phtateru
\[\lambda=\frac{1}{\tilde\nu}\]
\phtoku
\[\lambda=\frac{1}{2885}\ \text{cm}\approx\qty{3.47e-4}{cm}=\qty{3.47}{\micro m}\]
\phtashikame
可視光（\qtyrange{0.4}{0.78}{\micro m}）より長い赤外線です。H は Cl よりずっと軽いので、換算質量が小さく、CO より振動数が高く、波長が短くなっています。
\kotae 約 \qty{3.47}{\micro m}
\end{phreidai}
\begin{jibun}
CO（\qty{2143}{cm^{-1}}）が吸収する赤外線の波長を求めましょう。
\bkjans{$1/2143$ cm で、約 \qty{4.67}{\micro m}}
\end{jibun}

\label{ph:原-6:ex10}
\begin{phreidai}[CO の零点のゆれの幅]
CO の振動のいちばん低い状態で、原子の間の距離がゆれる幅の目安 $x_0=\sqrt{\hbar/(\mu\omega)}$ を求め、結合の長さ \qty{112.8}{pm}と比べましょう。
\phbamen
零点エネルギーのため、いちばん低い状態でも分子がゆれていることを、大きさで見る場面です（⑧）。
\phmushi
調和振動子とします。
\phzu
$\hbar=\qty{1.055e-34}{J.s}$、$\mu=\qty{1.14e-26}{kg}$、$\omega=\qty{4.04e14}{rad/s}$。
\phriyu
$\psi_0=e^{-x^2/(2x_0^2)}$ の山の幅が $x_0$（⑧）。
\phtateru
\[x_0=\sqrt{\frac{\hbar}{\mu\omega}}\]
\phtoku
\[x_0=\sqrt{\frac{1.055\times10^{-34}}{1.14\times10^{-26}\times4.04\times10^{14}}}\approx\qty{4.8e-12}{m}=\qty{4.8}{pm}\]
\phtashikame
結合の長さの約4\,\% です。小さいけれど、0ではありません。どんなに冷やしても、分子はこのくらいゆれています。
\kotae 約 \qty{4.8}{pm}（結合の長さの約4\,\%）
\end{phreidai}
\begin{jibun}
HCl（$\mu=\qty{1.63e-27}{kg}$、$\omega=\qty{5.44e14}{rad/s}$）の $x_0$ を求め、結合の長さ \qty{127.5}{pm}と比べましょう。
\bkjans{$x_0\approx\qty{10.9}{pm}$。結合の長さの約9\,\%（H が軽いので大きくゆれる）}
\end{jibun}

\label{ph:原-6:ex11}
\begin{phreidai}[アンモニアのトンネル分裂]
アンモニアの分子の窒素が、3つの水素の面を通りぬけて行き来するトンネル分裂の振動数は \qty{23.87}{GHz}です。分裂のエネルギーを eV で求め、\qty{300}{K}の $k_{\mathrm B}T$ と比べましょう。
\phbamen
2つの井戸のトンネル分裂を、実物の分子で見る場面です（⑪）。
\phmushi
分子の回転による違い（線の分かれ）は考えません。
\phzu
$f=\qty{23.87e9}{Hz}$、$h=\qty{4.14e-15}{eV.s}$、$k_{\mathrm B}T=\qty{0.0259}{eV}$。
\phriyu
分裂のエネルギーは、上と下の準位の間をうつるときに出入りする光子のエネルギー $hf$。
\phtateru
\[\Delta E=hf\]
\phtoku
\[\Delta E=4.14\times10^{-15}\times2.387\times10^{10}\approx\qty{9.9e-5}{eV}\]
$k_{\mathrm B}T$ の約 $\tfrac{1}{260}$。上と下の準位の分子の数の比は $e^{-\Delta E/(k_{\mathrm B}T)}\approx0.996$。
\phtashikame
上と下の分子の数は、ほとんど同じです。だから、そのままでは誘導放出（原-4 ⑪）で光を強められません。メーザーは、電場で上の準位の分子だけをより分けて、この問題を解きました。
\kotae 約 $\qty{9.9e-5}{eV}$（$k_{\mathrm B}T$ よりずっと小さい）
\end{phreidai}
\begin{jibun}
\qty{23.87}{GHz}のマイクロ波の波長と、窒素が左右を行き来する周期 $1/f$ を求めましょう。
\bkjans{波長は $c/f$ で約 \qty{1.26}{cm}。周期は $1/f$ で約 $\qty{4.2e-11}{s}$（\qty{42}{ps}）}
\end{jibun}

\label{ph:原-6:ex12}
\begin{phreidai}[表計算で井戸の準位をさがす]
例題5の井戸（深さ \qty{1}{eV}・幅 \qty{1}{nm}）の準位を、表計算の「撃って合わせる」方法（原-5 ⑫）で求めましょう。
\phbamen
グラフの交点を使わずに、方程式を数値で進めて準位をさがす場面です。
\phmushi
2回微分を、となりの点の差で近づけます。
\phzu
まん中から右へ、$\Delta x=\qty{0.02}{nm}$ ごとの点 $x=0.01,\ 0.03,\ \ldots$ を並べる。井戸の中（$x<0.5$）は $U=0$、外は $U=1$。$\hbar^2/(2m)=\qty{0.0381}{eV.nm^2}$。
\phriyu
原-5 ⑫の差分の式で、$\psi$ を右へ順に進められる。正しい $E$ のときだけ、遠くで $\psi$ が0に近づき、少しずれると、正か負の大きな値に飛んでいく。
\phtateru
偶の解は、$x=-0.01$ と $x=0.01$ の値を1と1、奇の解は $-1$ と1にして、
\[\psi_{j+1}=2\psi_j-\psi_{j-1}+\frac{(0.02)^2}{0.0381}(U_j-E)\,\psi_j\]
で $x=\qty{2}{nm}$ まで進め、そこで符号が変わる $E$ をさがす。
\phtoku
偶：$E=0.19$ で $x=2$ の値は $+1.5$、$E=0.20$ で $-22$。奇：$E=0.70$ で $+9.1$、$E=0.71$ で $-28$。範囲をせばめていくと、\qty{0.191}{eV}と \qty{0.702}{eV}。
\phtashikame
例題5のグラフの交点（0.19・\qty{0.70}{eV}）と同じです。$E$ をセル1つで変えられるようにしておくと、表全体が一度に計算し直されます。
\kotae 約 \qty{0.19}{eV}と約 \qty{0.70}{eV}（例題5と同じ）
\end{phreidai}
\begin{jibun}
同じ表で、井戸の幅を \qty{2}{nm}（$x<1$ が井戸）にし、$x=\qty{3}{nm}$ まで進めて、いちばん低い準位（偶）をさがしましょう。$E=0.065$ と 0.066 で $x=3$ の値の符号はどうなりますか。
\bkjans{0.065 で正、0.066 で負なので、約 \qty{0.066}{eV}}
\end{jibun}

%% ======================== まちがえやすいところ ========================
\Needspace{11\baselineskip}\bksection{tsumazuki}{まちがえやすいところをチェック}
\begin{tsumazuki}{壁の厚さ2倍で、割合は半分？}
\bkng{厚さが2倍なら、すりぬける割合は $\tfrac12$。}
\bkok{$e^{-2\kappa a}$ の肩が2倍になるので、割合は2乗（例：$10^{-5}$ が $10^{-10}$）。}
指数関数では、「たす」が「かける」になります（③・5-14）。
\end{tsumazuki}

\begin{tsumazuki}{すりぬけると、エネルギーが減る？}
\bkng{壁をぬけた電子は、エネルギーを使ったので、エネルギーが小さくなっている。}
\bkok{壁の左右の $U$ が同じなら、ぬけた電子のエネルギーは来たときと同じ $E$。}
小さくなるのは、そこで見つかる確率（波の大きさ）です。エネルギー（波長）は変わりません（図\ref{ph:原-6:fig-shimi}）。
\end{tsumazuki}

\begin{tsumazuki}{壁より高いエネルギーなら、全部通る？}
\bkng{$E>U$ なら、ふつうの粒と同じく全部が段を上がる。}
\bkok{波は、段で一部がはね返る（$E=2U$ で約3\,\%）。}
波数が変わる所では、ひもの波と同じく反射が起きます（②）。
\end{tsumazuki}

\begin{tsumazuki}{調和振動子の準位も、上ほど間かくが広い？}
\bkng{箱の電子と同じく、準位は $n^2$ に比例し、上ほど間かくが広い。}
\bkok{調和振動子の準位は $(n+\tfrac12)\hbar\omega$ で、間かくはどこも $\hbar\omega$。}
井戸の形で、準位の並び方が変わります（原-5 の箱は $n^2$）。
\end{tsumazuki}

\begin{tsumazuki}{いちばん低い振動のエネルギーは0？}
\bkng{いちばん低い状態では、分子は止まっている。}
\bkok{いちばん低い状態にも $\tfrac12\hbar\omega$ が残り、分子はゆれている（零点エネルギー）。}
冷やしても0にならないので、比熱には効きませんが、分子の大きさや結合のエネルギーには効きます（⑧）。
\end{tsumazuki}

%% ======================== 工学ではこう使う ========================
\Needspace{10\baselineskip}\bksection{phgear}{工学ではこう使う}
\begin{phkougaku}[壁の厚さを nm で設計する]
\textbf{フラッシュメモリ・SSD}：記録を何年も保つには、酸化膜を厚くしたい。でも、厚すぎると書きこみに高い電圧と時間がかかります。そのつり合いで、酸化膜の厚さが決められています（④）。書きかえを何万回も繰り返すと、膜に小さな傷ができて電子がもれやすくなるので、SSD は書きこむ場所をまんべんなく使うように制御しています。

\textbf{磁気トンネル接合}：厚さ約 \qty{1}{nm}の絶縁膜を、2枚の磁石の膜ではさむと、2枚の磁石の向きがそろっているときだけ、トンネル電流がよく流れます。ハードディスクの読み取りや、電源を切っても忘れないメモリ（MRAM）に使われています。

\textbf{赤外分光の装置}：工場の排ガスや、飲酒の検査の呼気の成分を、分子のゆれが吸収する赤外線の波長と強さではかります（⑨）。

\begin{phketa}[トンネル効果と分子のゆれの数の目安]
\phketaline{壁の差 \qty{1}{eV}の $\kappa$}{約 \qty{5}{nm^{-1}}（$1/\kappa\approx\qty{0.2}{nm}$）}
\phketaline{STM：電流が10倍になる近づき}{約 \qty{0.1}{nm}}
\phketaline{フラッシュメモリの酸化膜}{約 \qty{8}{nm}・壁 約 \qty{3}{eV}}
\phketaline{分子の振動}{0.1〜\qty{0.5}{eV}（\qtyrange{2.5}{12}{\micro m}）}
\phboxhead{答えが桁で合っているか}
すりぬける割合は、肩の $2\kappa a$ を先に計算し、「2.3 で1けた」（$e^{2.3}\approx10$）で桁を読むと速く確かめられます。$2\kappa a=46$ なら約 $10^{-20}$ です。
\end{phketa}
\end{phkougaku}
\phstation{半導体デバイス,表面科学,分子分光学,超伝導エレクトロニクス}{トンネル効果、トンネル電流、浮遊ゲート、ファウラー・ノルドハイム、走査トンネル顕微鏡、零点エネルギー、赤外分光、波数（cm$^{-1}$）、WKB 近似}

%% ======================== 章末問題 ========================
\bkendproblems
\bkproblemgroup{基本}{3}
\bkproblem 高さの差 \qty{3}{eV}の壁について、（1）$\kappa$ を求めましょう。（2）厚さ \qty{1}{nm}と \qty{2}{nm}のすりぬける割合 $e^{-2\kappa a}$ を求めましょう。
\bkpans{（1）$5.12\times\sqrt3\approx\qty{8.87}{nm^{-1}}$。（2）\qty{1}{nm}：約 $\num{2.0e-8}$、\qty{2}{nm}：約 $\num{3.9e-16}$（2乗）}
\bkproblem 高さの差 \qty{1}{eV}の壁で、すりぬける割合 $e^{-2\kappa a}$ を $10^{-10}$ にするには、厚さを何 nm にすればよいですか（$\ln10^{10}\approx23.0$）。差が \qty{4}{eV}ならどうですか。
\bkpans{\qty{1}{eV}：$a=23.0\div(2\times5.12)\approx\qty{2.2}{nm}$。\qty{4}{eV}：$\kappa$ が2倍なので約 \qty{1.1}{nm}}
\bkproblem エネルギー \qty{4}{eV}の電子が、高さ \qty{3}{eV}の段差に来ました。反射率を求めましょう。
\bkpans{$k_2/k_1=\sqrt{1/4}=0.5$。$\mathcal R=(0.5/1.5)^2=\tfrac19\approx0.11$}
\bkproblem 例題1のフラッシュメモリの酸化膜（壁の高さとの差 \qty{3.2}{eV}）を、\qty{8}{nm}から \qty{7}{nm}に \qty{1}{nm}だけうすくすると、すりぬける割合 $e^{-2\kappa a}$ は何倍になりますか。
\bkpans{$\kappa\approx\qty{9.17}{nm^{-1}}$。割合は $e^{2\kappa\times1}=e^{18.3}\approx\num{9e7}$ 倍（約1億倍）}
\bkproblem 高さの差 \qty{1}{eV}・厚さ \qty{0.1}{nm}の壁を、電子と陽子（質量は電子の1836倍）がすりぬける割合 $e^{-2\kappa a}$ を、それぞれ求めましょう。
\bkpans{$\kappa\propto\sqrt m$。電子：$\kappa=\qty{5.12}{nm^{-1}}$、$e^{-1.02}\approx0.36$。陽子：$\kappa=5.12\times\sqrt{1836}\approx\qty{219}{nm^{-1}}$、$e^{-43.9}\approx\num{9e-20}$（重い粒ほどすりぬけにくい）}
\bkproblem 例題4の STM で、針を \qty{0.05}{nm}遠ざけると、トンネル電流は何倍になりますか。
\bkpans{$e^{-2\kappa\times0.05}=e^{-1.09}\approx0.34$ 倍（約3分の1）}
\bkproblemgroup{チャレンジ}{5}
\bkproblem 深さ \qty{2}{eV}、幅 \qty{0.5}{nm}の井戸について、（1）$z_0$ を求め、閉じこめられた状態の数を答えましょう。（2）いちばん低い準位を、グラフの交点か二分法で求めましょう。
\bkpans{（1）$z_0\approx1.81$ で、2つ。（2）偶の交点 $z\approx0.992$ から約 \qty{0.60}{eV}}
\bkproblem 調和振動子の $\psi_1=xe^{-x^2/(2x_0^2)}$ を、シュレディンガー方程式に代入して、エネルギーが $\tfrac32\hbar\omega$ であることを確かめましょう。
\bkpans{$\psi_1''=\Bigl(\dfrac{x^2}{x_0^{\,4}}-\dfrac{3}{x_0^{\,2}}\Bigr)\psi_1$。代入すると $x^2$ の項は⑧と同じく打ち消し合い、$\dfrac{3\hbar^2}{2mx_0^{\,2}}\psi_1=\dfrac32\hbar\omega\,\psi_1$ が残る}
\bkproblem 窒素（N$_2$）の振動は \qty{2330}{cm^{-1}}です。（1）$\hbar\omega$ と $\theta_{\mathrm E}$、（2）\qty{300}{K}の $C_{\text{振}}/R$ を求め、空気の定積モル比熱が約 $\tfrac52R$ のわけを説明しましょう（$e^{11.2}\approx\num{7.3e4}$）。
\bkpans{（1）約 \qty{0.289}{eV}、約 \qty{3350}{K}。（2）約0.0017。振動はこおり、進む動き3つと回る動き2つの $\tfrac52R$ だけになる}
\bkproblem ⑪のトンネル分裂は、壁の厚さ $b$ のおよそ $e^{-\kappa b}$ に比例します。壁との差 \qty{0.81}{eV}の $\kappa$ で、壁を \qty{0.5}{nm}から \qty{1.0}{nm}にすると分裂が何倍になるか見積もり、差分法の値と比べましょう。
\bkpans{$\kappa=5.12\times\sqrt{0.81}\approx\qty{4.61}{nm^{-1}}$。$e^{-4.61\times0.5}\approx0.10$ 倍で、差分法の $1.87\div18.7=0.10$ と合う}

\begin{furikaeri}
\begin{checklist}
\item 壁の中で波動関数が $e^{-\kappa x}$ で弱まることを使い、すりぬける割合 $e^{-2\kappa a}$ を見積もれる。\bkfurref{ph:原-6:sec3}
\item フラッシュメモリや走査トンネル顕微鏡のしくみを、トンネル効果で説明できる。\bkfurref{ph:原-6:sec4}
\item 境目のつながりから、階段の反射率・有限の井戸の準位・正確な透過率を求められる。\bkfurref{ph:原-6:sec6}
\item 調和振動子の準位 $(n+\tfrac12)\hbar\omega$ を使い、分子の振動の光と比熱を説明できる。\bkfurref{ph:原-6:sec8}
\end{checklist}
\end{furikaeri}
\bkfukidashi{mina}{USBメモリの写真が消えないのは、電子がガラスの壁をすりぬけられないから。書きこむときは、そのすりぬけを使うなんて、おもしろい！}
